\documentclass[11pt]{article}
% ============================================================================
% Supplement to the JASA-structured MWQTM manuscript.
%   A. Proofs            (transport/identification/orthogonality/remainder/
%                         one-step normality/sandwich; clustered Bahadur;
%                         censoring orthogonality)  -- imported proof bodies
%   B. Algorithms and implementation
%   C. Additional simulations (audit and sensitivity tables)
%   D. Additional application details
% ============================================================================
\usepackage[margin=1in]{geometry}
\usepackage{amsmath,amssymb,amsthm,bm,booktabs,graphicx,natbib,hyperref}
\usepackage{xr}
\externaldocument{paper}                      % cross-reference the main paper
\hypersetup{colorlinks=true,linkcolor=blue,citecolor=blue,urlcolor=blue}
\newtheorem{theorem}{Theorem}
\newtheorem{proposition}{Proposition}
\newtheorem{lemma}{Lemma}
\newtheorem{corollary}{Corollary}
\newtheorem{assumption}{Assumption}
\newtheorem{condition}{Condition}
\theoremstyle{remark}\newtheorem{remark}{Remark}
\newcommand{\E}{\mathbb{E}}
\newcommand{\PR}{\mathbb{P}}
\newcommand{\Pn}{\mathbb{P}_n}
\newcommand{\Gn}{\mathbb{G}_n}
\newcommand{\X}{\bm{X}}
\newcommand{\Z}{\bm{Z}}
\newcommand{\bbeta}{\bm{\beta}}
\newcommand{\bgamma}{\bm{\gamma}}
\newcommand{\bth}{\bm{\theta}}
\newcommand{\indc}[1]{\mathbf{1}\{#1\}}
\newcommand{\given}{\,\vert\,}
\newcommand{\binit}{\widehat\bbeta_{\rm init}}
\newcommand{\rom}{r_{\omega}}
\newcommand{\op}{\mathrm{op}}
\DeclareMathOperator*{\argmin}{arg\,min}
\DeclareMathOperator*{\argmax}{arg\,max}

\renewcommand{\thesection}{S\arabic{section}}
\renewcommand{\thetable}{S\arabic{table}}
\renewcommand{\thefigure}{S\arabic{figure}}
\renewcommand{\theequation}{S\arabic{equation}}

\title{\bf Supplement to ``Event-aligned conditional quantile trajectories of a
biomarker before disease onset under left truncation and competing death''}
\author{\thanks{}}
\date{\today}

\begin{document}
\maketitle
\tableofcontents
\bigskip

This supplement collects the proofs (Section~\ref{supp:proofs}), the estimation
algorithms and implementation details (Section~\ref{supp:algo}), the full set of
simulation and audit/sensitivity studies summarized in Section~\ref{sec:sim} of the
main paper (Section~\ref{supp:sims}), and additional details of the application
(Section~\ref{supp:app}). Equations, tables, and figures are numbered S1, S2, \dots;
references of the form ``Theorem~1'' point to the main paper.

% ============================================================================
\section{Proofs}\label{supp:proofs}
This section reproduces the asymptotic theory: the transport/identification
arguments, the implicit density-weighted orthogonalization, the censoring
orthogonality ($\chi^C=0$), the sharper-remainder bound, the high-level
cross-fitted one-step asymptotic normality and sandwich consistency, and the
clustered, weighted, growing-sieve Gram/Knight/Bahadur representation that verifies
the high-level conditions for the concrete sieve. The numbering of internal
results is self-contained within this section.

\subsection*{S1.A\quad Identification, orthogonality, one-step normality, sandwich}
\input{proof_body}

\subsection*{S1.B\quad Censoring orthogonality of the IPW quantile score}
\input{censoring_orthogonality_body}

\subsection*{S1.C\quad Clustered weighted growing-sieve Bahadur representation}
\input{bahadur_body}

% ============================================================================
\section{Algorithms and implementation}\label{supp:algo}

\paragraph{Sieve basis.} The baseline surface $m_{\tau}(u,t)$ uses a tensor-product
cubic B-spline basis $\{B_k(u,t)\}$ over the two time scales, a partition of unity
(so no separate intercept enters the design). The implementation uses this
partition-of-unity basis; the asymptotic analysis (Section~\ref{supp:proofs}) uses
the $L_2$-normalized basis $\widetilde B_k=\sqrt K\,B_k$, under which the population
Gram has eigenvalues bounded away from $0$ and $\infty$; the two differ by the
fixed scaling $\sqrt K$, which cancels in the estimator.

\paragraph{Complexity selection (A-CV).} The number of interior knots / spline
degrees of freedom and any penalty $\lambda_n$ are chosen by subject-level
$K$-fold cross-validation minimizing the held-out check loss (the ``-CV'' suffix).
A fixed richer-sieve fit (``A-Rich'', $K{+}2$) is reported as a sensitivity
comparison, not a theory-guided undersmoothing rule. The reported pilot uses a
regression-spline sieve (fixed knots, $\lambda_n=0$).

\paragraph{Sparsity and projection (B-OS-CF).} The conditional density at the
quantile, $f_\tau(0)$, is estimated by the Hendricks--Koenker quantile-spacing
sparsity $\widehat f_\tau(0)=\{(\widehat q_{\tau+a_n}-\widehat q_{\tau-a_n})/(2a_n)\}^{-1}$,
rearranged for positivity and truncated to $[f_{\min},f_{\max}]$. The
density-weighted projection $\widehat g_{\tau,\omega}$ solves a weighted spline
least-squares regression of $\X$ on $B_g(u,t)$ with weights
$a_\omega\widehat f_\tau(0)$. Cross-fitting over subject folds is part of the
estimator definition (Section~\ref{sec:onestep} of the main paper).

\paragraph{Bootstrap and fold discipline.} Inference for the A-CV point estimator
uses a full-pipeline subject-level bootstrap: subjects (not observations) are
resampled, preserving within-subject clustering, and the entire pipeline
(complexity selection and, for B-OS-CF, nuisance fitting) is refit on each
resample. Cross-validation folds are formed by subject id, so no subject's visits
appear in both training and held-out folds; a controlled same-dataset comparison
(Table~\ref{tab:leakfold}) quantifies the effect of violating this fold discipline.

\paragraph{Crossing.} Predicted quantile curves are made monotone in $\tau$ by
rearrangement; this is a pointwise post-processing step that does not alter the raw
$\widehat\bbeta_\tau$ used for inference.

\paragraph{Reproducibility and computational audit.} A computational-failure audit
(Table~\ref{tab:audit}) records the disposition of every replicate (QR-solver
status, degenerate folds, support-mask exclusions) so that reported summaries are
not silently conditioned on ``clean'' replicates. Code is available at the
companion repository.

% ============================================================================
\section{Additional simulations}\label{supp:sims}
This section gives the full designs and complete tables for the studies summarized
in Section~\ref{sec:sim} of the main paper. Unless noted, $R$ denotes the number of
Monte-Carlo replicates and coverage is for the $95\%$ subject-bootstrap interval.

\subsection{Estimator family and finite-sample inference}
\begin{table}[t]\centering\small
\caption{Headline estimator comparison on a \emph{single} paired experiment ($R=1000$; bootstrap coverage over a $500$-replicate subset, $B=500$ full-refit resamples; MC SE $\approx0.010$). Four estimators on the same replicates: complete-event (CC, $\omega=1$) and IPCW ($\omega=1/\widehat G_C$), each as the joint-sieve A-CV point estimator and the cross-fitted one-step \textbf{B-OS-CF} ($K=5$ folds). All coverages are full-refit subject bootstrap.}
\label{tab:estsim}\begin{tabular}{rrr r rrrr rrrr}\toprule
&&&& \multicolumn{4}{c}{bias} & \multicolumn{4}{c}{bootstrap coverage}\\
\cmidrule(lr){5-8}\cmidrule(lr){9-12}
Sc.&$n$&$\tau$&$\beta_{\tau,1}$& \tiny CC-ACV & \tiny IP-ACV & \tiny CC-CF & \tiny IP-CF & \tiny CC-ACV & \tiny IP-ACV & \tiny CC-CF & \tiny IP-CF\\\midrule
1 & 500 & 0.25 & 0.500 & -0.002 & -0.002 & 0.000 & -0.000 & 0.960 & 0.958 & 0.970 & 0.968 \\
1 & 500 & 0.50 & 0.500 & -0.000 & 0.001 & 0.002 & 0.002 & 0.968 & 0.962 & 0.984 & 0.968 \\
1 & 500 & 0.75 & 0.500 & -0.002 & -0.003 & -0.008 & -0.009 & 0.950 & 0.954 & 0.966 & 0.962 \\
\addlinespace
1 & 1000 & 0.25 & 0.500 & 0.004 & 0.004 & 0.004 & 0.005 & 0.938 & 0.942 & 0.960 & 0.960 \\
1 & 1000 & 0.50 & 0.500 & -0.001 & -0.001 & -0.001 & -0.001 & 0.946 & 0.948 & 0.954 & 0.960 \\
1 & 1000 & 0.75 & 0.500 & 0.002 & 0.002 & 0.001 & 0.003 & 0.958 & 0.946 & 0.958 & 0.972 \\
\addlinespace
1 & 2000 & 0.25 & 0.500 & 0.000 & -0.000 & 0.001 & 0.001 & 0.956 & 0.952 & 0.970 & 0.966 \\
1 & 2000 & 0.50 & 0.500 & -0.000 & -0.000 & 0.000 & -0.000 & 0.966 & 0.956 & 0.966 & 0.978 \\
1 & 2000 & 0.75 & 0.500 & 0.002 & 0.002 & 0.002 & 0.003 & 0.942 & 0.944 & 0.972 & 0.966 \\
\addlinespace
2 & 500 & 0.25 & 0.230 & 0.006 & 0.007 & 0.004 & 0.002 & 0.918 & 0.928 & 0.954 & 0.952 \\
2 & 500 & 0.50 & 0.500 & -0.002 & -0.003 & -0.004 & -0.004 & 0.950 & 0.954 & 0.948 & 0.948 \\
2 & 500 & 0.75 & 0.770 & -0.005 & -0.005 & -0.002 & -0.003 & 0.954 & 0.958 & 0.954 & 0.958 \\
\addlinespace
2 & 1000 & 0.25 & 0.230 & 0.001 & 0.003 & -0.001 & -0.001 & 0.940 & 0.932 & 0.970 & 0.974 \\
2 & 1000 & 0.50 & 0.500 & 0.002 & 0.003 & 0.002 & 0.003 & 0.946 & 0.938 & 0.972 & 0.962 \\
2 & 1000 & 0.75 & 0.770 & -0.005 & -0.007 & -0.004 & -0.005 & 0.946 & 0.946 & 0.958 & 0.970 \\
\addlinespace
2 & 2000 & 0.25 & 0.230 & -0.001 & -0.001 & -0.002 & -0.001 & 0.938 & 0.940 & 0.960 & 0.956 \\
2 & 2000 & 0.50 & 0.500 & -0.000 & 0.000 & -0.001 & -0.000 & 0.944 & 0.954 & 0.964 & 0.956 \\
2 & 2000 & 0.75 & 0.770 & 0.002 & 0.002 & 0.003 & 0.004 & 0.948 & 0.950 & 0.980 & 0.986 \\
\bottomrule\end{tabular}\end{table}

\begin{table}[t]\centering\small
\caption{Reporting-completeness supplement to Tables~\ref{tab:estsim}/\ref{tab:estse}: the Monte-Carlo ESD and mean bootstrap SE ($\overline{\mathrm{SE}}$) for all four estimators, reported \emph{directly} so neither need be backed out of SER/RMSE ($R=1000$, $R_{\rm cov}=500$). Within each pipeline B-OS-CF has a larger ESD than A-CV. The B-OS-CF$-$A-CV \emph{difference} is dominated in \emph{magnitude} by the orthogonal correction, with a smaller sample-splitting contribution (Table~\ref{tab:corr}: $|c_{\rm os}|/$ESD $\approx0.28$--$0.40$ vs $|c_{\rm split}|/$ESD $\approx0.07$--$0.13$); we do \emph{not} decompose the variance of the ESD increase itself (that would require the covariance terms with $\widehat\beta_{\rm ACV}$). that $\overline{\mathrm{SE}}\gtrsim$ESD for B-OS-CF is consistent with bootstrap-SE overestimation being an important contributor to its mild over-coverage. Bias and coverage are in Table~\ref{tab:estsim}; SER and RMSE in Table~\ref{tab:estse}.}
\label{tab:estfull}\begin{tabular}{rrr rr rr rr rr}\toprule
&&&\multicolumn{2}{c}{CC A-CV}&\multicolumn{2}{c}{CC B-OS-CF}&\multicolumn{2}{c}{IP A-CV}&\multicolumn{2}{c}{IP B-OS-CF}\\
\cmidrule(lr){4-5}\cmidrule(lr){6-7}\cmidrule(lr){8-9}\cmidrule(lr){10-11}
Sc.&$n$&$\tau$&ESD&$\overline{\rm SE}$&ESD&$\overline{\rm SE}$&ESD&$\overline{\rm SE}$&ESD&$\overline{\rm SE}$\\\midrule
1 & 500 & 0.25 & 0.110 & 0.113 & 0.123 & 0.129 & 0.115 & 0.117 & 0.131 & 0.135 \\
1 & 500 & 0.50 & 0.099 & 0.108 & 0.106 & 0.121 & 0.106 & 0.112 & 0.117 & 0.126 \\
1 & 500 & 0.75 & 0.108 & 0.113 & 0.117 & 0.129 & 0.112 & 0.117 & 0.124 & 0.134 \\
\addlinespace
1 & 1000 & 0.25 & 0.078 & 0.079 & 0.087 & 0.094 & 0.083 & 0.083 & 0.092 & 0.100 \\
1 & 1000 & 0.50 & 0.073 & 0.075 & 0.082 & 0.088 & 0.077 & 0.079 & 0.086 & 0.093 \\
1 & 1000 & 0.75 & 0.079 & 0.080 & 0.089 & 0.095 & 0.083 & 0.083 & 0.095 & 0.100 \\
\addlinespace
1 & 2000 & 0.25 & 0.055 & 0.056 & 0.059 & 0.067 & 0.059 & 0.059 & 0.063 & 0.071 \\
1 & 2000 & 0.50 & 0.049 & 0.053 & 0.054 & 0.061 & 0.052 & 0.055 & 0.058 & 0.065 \\
1 & 2000 & 0.75 & 0.055 & 0.056 & 0.059 & 0.065 & 0.058 & 0.058 & 0.065 & 0.070 \\
\addlinespace
2 & 500 & 0.25 & 0.144 & 0.145 & 0.160 & 0.167 & 0.152 & 0.151 & 0.168 & 0.175 \\
2 & 500 & 0.50 & 0.134 & 0.138 & 0.150 & 0.157 & 0.140 & 0.144 & 0.157 & 0.164 \\
2 & 500 & 0.75 & 0.142 & 0.146 & 0.160 & 0.169 & 0.148 & 0.152 & 0.165 & 0.176 \\
\addlinespace
2 & 1000 & 0.25 & 0.104 & 0.102 & 0.113 & 0.121 & 0.110 & 0.107 & 0.120 & 0.129 \\
2 & 1000 & 0.50 & 0.094 & 0.096 & 0.105 & 0.112 & 0.099 & 0.101 & 0.111 & 0.119 \\
2 & 1000 & 0.75 & 0.104 & 0.102 & 0.114 & 0.121 & 0.109 & 0.107 & 0.122 & 0.129 \\
\addlinespace
2 & 2000 & 0.25 & 0.073 & 0.072 & 0.078 & 0.084 & 0.077 & 0.075 & 0.084 & 0.090 \\
2 & 2000 & 0.50 & 0.068 & 0.068 & 0.073 & 0.078 & 0.072 & 0.071 & 0.078 & 0.084 \\
2 & 2000 & 0.75 & 0.071 & 0.072 & 0.075 & 0.084 & 0.074 & 0.076 & 0.080 & 0.090 \\
\bottomrule\end{tabular}\end{table}

\begin{table}[t]\centering\small
\caption{Supplement to Table~\ref{tab:estsim}: SER$=\overline{\mathrm{SE}}_{\rm boot}/\mathrm{ESD}$ and RMSE for \emph{all four} estimators ($R=1000$, $R_{\rm cov}=500$). Since bias is negligible in these cells, ESD$\approx$RMSE and $\overline{\mathrm{SE}}=\mathrm{SER}\times\mathrm{ESD}$; the underlying ESD and $\overline{\mathrm{SE}}$ are reported directly in Table~\ref{tab:estfull}. \emph{Most} coverage dips coincide with a SER slightly below $1$ (a mildly small bootstrap SE), \emph{not} bias; the exception is the single lowest-coverage cell (Sc.~2, $n{=}500$, $\tau{=}0.25$: cov $0.918$), which has SER $\approx1.00$ \emph{and} negligible bias ($0.006$), so its shortfall is neither mean SE nor bias. The studentization diagnostic (Table~\ref{tab:studentized}) shows $Z=(\widehat\beta-\beta_0)/\widehat{\mathrm{SE}}$ is near-Gaussian in its central moments (skew $-0.01$, excess kurtosis $0.02$) with normal $=$ percentile coverage, so the shortfall is not explained by mean SE, central skew/kurtosis, or a simple SE--error association; it is \emph{consistent with} a mild small-sample upper-tail deviation and shrinks with $n$. B-OS-CF carries a modestly larger RMSE, generally around $5$--$15\%$ (the variance cost of the explicit orthogonal adjustment), and---after the fold-leakage fix---a SER $\ge1$, mildly to \emph{moderately} conservative in several cells (coverage up to $\approx0.986$).}
\label{tab:estse}\begin{tabular}{rrr rr rr rr rr}\toprule
&&& \multicolumn{2}{c}{A-CV cc} & \multicolumn{2}{c}{A-CV ip} & \multicolumn{2}{c}{B-OS-CF cc} & \multicolumn{2}{c}{B-OS-CF ip}\\
\cmidrule(lr){4-5}\cmidrule(lr){6-7}\cmidrule(lr){8-9}\cmidrule(lr){10-11}
Sc.&$n$&$\tau$&SER&RMSE&SER&RMSE&SER&RMSE&SER&RMSE\\\midrule
1 & 500 & 0.25 & 1.03 & 0.110 & 1.02 & 0.115 & 1.05 & 0.123 & 1.03 & 0.131 \\
1 & 500 & 0.50 & 1.08 & 0.099 & 1.06 & 0.106 & 1.14 & 0.106 & 1.08 & 0.117 \\
1 & 500 & 0.75 & 1.05 & 0.108 & 1.05 & 0.112 & 1.10 & 0.117 & 1.09 & 0.124 \\
\addlinespace
1 & 1000 & 0.25 & 1.01 & 0.078 & 1.00 & 0.083 & 1.09 & 0.087 & 1.09 & 0.092 \\
1 & 1000 & 0.50 & 1.03 & 0.073 & 1.02 & 0.077 & 1.07 & 0.082 & 1.08 & 0.086 \\
1 & 1000 & 0.75 & 1.01 & 0.079 & 1.00 & 0.083 & 1.06 & 0.089 & 1.05 & 0.095 \\
\addlinespace
1 & 2000 & 0.25 & 1.01 & 0.055 & 1.00 & 0.059 & 1.12 & 0.059 & 1.13 & 0.063 \\
1 & 2000 & 0.50 & 1.07 & 0.049 & 1.07 & 0.052 & 1.13 & 0.054 & 1.13 & 0.058 \\
1 & 2000 & 0.75 & 1.00 & 0.055 & 1.00 & 0.058 & 1.10 & 0.059 & 1.08 & 0.065 \\
\addlinespace
2 & 500 & 0.25 & 1.00 & 0.144 & 0.99 & 0.152 & 1.04 & 0.160 & 1.04 & 0.168 \\
2 & 500 & 0.50 & 1.03 & 0.134 & 1.02 & 0.140 & 1.05 & 0.150 & 1.05 & 0.157 \\
2 & 500 & 0.75 & 1.03 & 0.142 & 1.03 & 0.148 & 1.05 & 0.160 & 1.07 & 0.165 \\
\addlinespace
2 & 1000 & 0.25 & 0.98 & 0.104 & 0.97 & 0.110 & 1.08 & 0.113 & 1.08 & 0.120 \\
2 & 1000 & 0.50 & 1.02 & 0.094 & 1.01 & 0.099 & 1.07 & 0.105 & 1.07 & 0.111 \\
2 & 1000 & 0.75 & 0.98 & 0.104 & 0.98 & 0.110 & 1.06 & 0.114 & 1.06 & 0.122 \\
\addlinespace
2 & 2000 & 0.25 & 0.98 & 0.073 & 0.98 & 0.077 & 1.08 & 0.078 & 1.07 & 0.084 \\
2 & 2000 & 0.50 & 0.99 & 0.068 & 0.99 & 0.072 & 1.07 & 0.073 & 1.07 & 0.078 \\
2 & 2000 & 0.75 & 1.02 & 0.071 & 1.02 & 0.074 & 1.13 & 0.075 & 1.13 & 0.080 \\
\bottomrule\end{tabular}\end{table}

\begin{table}[t]\centering\small
\caption{Cross-fitted one-step \textbf{B-OS-CF} ($K=5$), CC and IPCW members. corr is the per-cell \emph{mean} signed one-step correction $\widetilde\beta_{\rm CF}-\widehat\beta_{\rm ACV}$; std is the mean of the replicate-wise standardized corrections $\frac1R\sum_r\mathrm{corr}_r/\widehat{\mathrm{SE}}_r$ (so $\mathrm{corr}$ and $\mathrm{std}$ can differ in sign). Both signed summaries are small ($|\mathrm{corr}|\le0.007$, $|\mathrm{std}|\le0.05$ in every cell), consistent with the A-CV first-order conditions already nearly annihilating the orthogonal score at $\lambda=0$. \emph{Caveat}: a signed mean near zero can mask per-replicate cancellation; the per-replicate \emph{absolute} correction and its split/one-step decomposition are in Table~\ref{tab:corr}.}
\label{tab:bosc}\begin{tabular}{rrr rr rr}\toprule
Sc.&$n$&$\tau$&corr$_{\rm CC}$&corr$_{\rm IPCW}$&std$_{\rm CC}$&std$_{\rm IPCW}$\\\midrule
1 & 500 & 0.25 & 0.002 & 0.001 & 0.025 & 0.010 \\
1 & 500 & 0.50 & 0.002 & 0.002 & 0.019 & 0.003 \\
1 & 500 & 0.75 & -0.006 & -0.005 & -0.047 & -0.044 \\
\addlinespace
1 & 1000 & 0.25 & -0.000 & 0.001 & 0.006 & 0.007 \\
1 & 1000 & 0.50 & 0.000 & -0.001 & 0.009 & 0.005 \\
1 & 1000 & 0.75 & -0.001 & 0.001 & -0.019 & -0.004 \\
\addlinespace
1 & 2000 & 0.25 & 0.001 & 0.001 & 0.021 & 0.025 \\
1 & 2000 & 0.50 & 0.000 & -0.000 & 0.008 & -0.001 \\
1 & 2000 & 0.75 & -0.000 & 0.000 & -0.015 & 0.004 \\
\addlinespace
2 & 500 & 0.25 & -0.003 & -0.005 & -0.024 & -0.032 \\
2 & 500 & 0.50 & -0.002 & -0.001 & -0.023 & -0.004 \\
2 & 500 & 0.75 & 0.003 & 0.002 & 0.013 & -0.008 \\
\addlinespace
2 & 1000 & 0.25 & -0.003 & -0.004 & -0.006 & -0.044 \\
2 & 1000 & 0.50 & 0.000 & 0.000 & 0.019 & 0.011 \\
2 & 1000 & 0.75 & 0.001 & 0.001 & 0.002 & 0.009 \\
\addlinespace
2 & 2000 & 0.25 & -0.000 & 0.000 & 0.003 & 0.008 \\
2 & 2000 & 0.50 & -0.000 & -0.000 & -0.009 & -0.008 \\
2 & 2000 & 0.75 & 0.001 & 0.001 & 0.015 & 0.012 \\
\bottomrule\end{tabular}\end{table}


\subsection{Censoring weights and orthogonality}
The complete-event (CC), IPCW, and oracle-weight (ORC) estimators, the
estimated-vs-oracle weight-gap rate, and hard-censoring behavior.
\begin{table}[t]\centering\small
\caption{\textbf{Paired} CC--A-CV vs.\ IPCW--A-CV on the same replicates and bootstrap draws ($R=1000$; coverage over $500$; $B=500$). $\Delta$cov $=\mathrm{cov}_{\rm CC}-\mathrm{cov}_{\rm IPCW}$; McNemar $(b,c)$ are the discordant counts (CC covers \& IPCW not; IPCW covers \& CC not) with the two-sided \emph{exact} (binomial) McNemar $p$; ESD ratio is CC/IPCW. Of the $18$ cells, \emph{one} (Sc.~1, $n{=}1000$, $\tau{=}0.75$: $b{=}6,c{=}0$) gives an unadjusted exact McNemar $p=0.031$; \textbf{no cell remains significant after a Holm/Bonferroni multiplicity correction} (the next smallest is $p=0.07$). The absolute coverage differences are $|\Delta\mathrm{cov}|\le0.012$ with no systematic direction. At $R_{\rm cov}=500$ this does not formally establish equivalence at a $0.02$ margin (a paired TOST would, and is left to a larger run) but is consistent with it. Separately, the CC ESD is smaller in every cell---an \emph{observed} efficiency improvement under these DGPs (not a claim of provably greater accuracy, which a non-significant test cannot support).}
\label{tab:cc}\begin{tabular}{rrr rr r rr r r}\toprule
Sc.&$n$&$\tau$&cov$_{\rm CC}$&cov$_{\rm IPCW}$&$\Delta$cov&$b$&$c$&$p$&$\tfrac{\rm ESD_{CC}}{\rm ESD_{IPCW}}$\\\midrule
1 & 500 & 0.25 & 0.960 & 0.958 & 0.002 & 3 & 2 & 1.00 & 0.955 \\
1 & 500 & 0.50 & 0.968 & 0.962 & 0.006 & 6 & 3 & 0.51 & 0.936 \\
1 & 500 & 0.75 & 0.950 & 0.954 & -0.004 & 2 & 4 & 0.69 & 0.959 \\
\addlinespace
1 & 1000 & 0.25 & 0.938 & 0.942 & -0.004 & 3 & 5 & 0.73 & 0.942 \\
1 & 1000 & 0.50 & 0.946 & 0.948 & -0.002 & 3 & 4 & 1.00 & 0.950 \\
1 & 1000 & 0.75 & 0.958 & 0.946 & 0.012 & 7 & 1 & 0.07 & 0.949 \\
\addlinespace
1 & 2000 & 0.25 & 0.956 & 0.952 & 0.004 & 4 & 2 & 0.69 & 0.942 \\
1 & 2000 & 0.50 & 0.966 & 0.956 & 0.010 & 6 & 1 & 0.13 & 0.953 \\
1 & 2000 & 0.75 & 0.942 & 0.944 & -0.002 & 7 & 8 & 1.00 & 0.956 \\
\addlinespace
2 & 500 & 0.25 & 0.918 & 0.928 & -0.010 & 3 & 8 & 0.23 & 0.951 \\
2 & 500 & 0.50 & 0.950 & 0.954 & -0.004 & 3 & 5 & 0.73 & 0.954 \\
2 & 500 & 0.75 & 0.954 & 0.958 & -0.004 & 2 & 4 & 0.69 & 0.958 \\
\addlinespace
2 & 1000 & 0.25 & 0.940 & 0.932 & 0.008 & 9 & 5 & 0.42 & 0.947 \\
2 & 1000 & 0.50 & 0.946 & 0.938 & 0.008 & 5 & 1 & 0.22 & 0.945 \\
2 & 1000 & 0.75 & 0.946 & 0.946 & 0.000 & 4 & 4 & 1.00 & 0.948 \\
\addlinespace
2 & 2000 & 0.25 & 0.938 & 0.940 & -0.002 & 2 & 3 & 1.00 & 0.949 \\
2 & 2000 & 0.50 & 0.944 & 0.954 & -0.010 & 1 & 6 & 0.13 & 0.951 \\
2 & 2000 & 0.75 & 0.948 & 0.950 & -0.002 & 3 & 4 & 1.00 & 0.950 \\
\bottomrule\end{tabular}\end{table}

\begin{table}[t]\centering\small
\caption{Identification/efficiency and the $\chi^C=0$ diagnostic (A-CV $\beta_{\tau,1}$, $\tau=0.5$, $R=400$; base-censoring regime). CC = unweighted complete-event; IPCW = estimated $\widehat G_C$; ORC = oracle $G_{C,0}$. All three are unbiased; ESD$_{\rm IPCW}\approx$ESD$_{\rm ORC}$ (no detectable additional variance from estimating $G_C$) and $\ge$ESD$_{\rm CC}$. The gap $\mathrm{sd}(\widehat\beta_{\rm IPCW}-\widehat\beta_{\rm ORC})/$ESD$_{\rm IPCW}$ shrinks with $n$ ($n^{-1}$-like, $\ll n^{-1/2}$).}
\label{tab:compare}\begin{tabular}{rr rrr rrr r}\toprule
Sc.&$n$&bias$_{\rm CC}$&bias$_{\rm IPCW}$&bias$_{\rm ORC}$&ESD$_{\rm CC}$&ESD$_{\rm IPCW}$&ESD$_{\rm ORC}$&gap/ESD\\\midrule
1 & 500 & -0.0033 & -0.0041 & -0.0034 & 0.1011 & 0.1057 & 0.1044 & 0.1163 \\
1 & 1000 & 0.0012 & 0.0010 & 0.0011 & 0.0728 & 0.0766 & 0.0762 & 0.0842 \\
1 & 2000 & -0.0010 & -0.0011 & -0.0012 & 0.0482 & 0.0501 & 0.0502 & 0.0710 \\
1 & 4000 & 0.0040 & 0.0045 & 0.0044 & 0.0360 & 0.0377 & 0.0377 & 0.0483 \\
\addlinespace
2 & 500 & -0.0049 & -0.0039 & -0.0041 & 0.1360 & 0.1460 & 0.1455 & 0.1206 \\
2 & 1000 & 0.0020 & 0.0027 & 0.0031 & 0.1007 & 0.1070 & 0.1065 & 0.0783 \\
2 & 2000 & -0.0015 & -0.0013 & -0.0011 & 0.0664 & 0.0702 & 0.0700 & 0.0621 \\
2 & 4000 & -0.0043 & -0.0048 & -0.0048 & 0.0511 & 0.0540 & 0.0539 & 0.0495 \\
\bottomrule\end{tabular}\end{table}

\begin{table}[t]\centering\small
\caption{Censoring-weight stress test: the $\chi^C=0$ gap under the base regime vs.\ a \emph{hard} regime (later/wider entry $+$ heavier, more covariate-dependent loss; Sc.\ 2, $\tau=0.5$, $R=400$). The hard regime lifts censoring to $\approx43\%$ with markedly heavier IPCW-weight tails ($99$th pct $\approx5$, max $\approx12$) and lower weight ESS$/n_T$ ($\approx0.79$ vs.\ $0.98$), i.e.\ closer to the positivity boundary; the Cox $\widehat G_C$ never failed. Even so, the estimated-vs-oracle gap/ESD stays small relative to the estimator ESD and decreases overall ($0.20,0.19,0.13,0.13$ at $n=500,1000,2000,4000$); we do \emph{not} claim a clean $n^{-1}$ rate---the final two sample sizes are nearly flat---only that the ratio stays small and decreases overall, which is \emph{consistent with} (but does not establish) the second-order, $\ll n^{-1/2}$ behaviour. ($\tau=0.5$ only; tail behaviour under hard weights is not probed.)}
\label{tab:censhard}\begin{tabular}{lr rr rrrr r}\toprule
regime&$n$&\%cens&ESS/$n_T$&$\widetilde w$&$w_{95}$&$w_{99}$&$w_{\max}$&gap/ESD\\\midrule
base & 500 & 0.17 & 0.98 & 1.17 & 1.53 & 1.76 & 2.07 & 0.1206 \\
base & 1000 & 0.17 & 0.98 & 1.17 & 1.53 & 1.75 & 2.14 & 0.0783 \\
base & 2000 & 0.17 & 0.98 & 1.17 & 1.53 & 1.74 & 2.19 & 0.0621 \\
base & 4000 & 0.17 & 0.98 & 1.17 & 1.52 & 1.73 & 2.29 & 0.0495 \\
\addlinespace
hard & 500 & 0.43 & 0.79 & 1.47 & 3.35 & 5.10 & 7.43 & 0.1997 \\
hard & 1000 & 0.42 & 0.79 & 1.47 & 3.40 & 5.16 & 8.48 & 0.1925 \\
hard & 2000 & 0.43 & 0.79 & 1.47 & 3.42 & 5.28 & 9.66 & 0.1308 \\
hard & 4000 & 0.43 & 0.78 & 1.47 & 3.42 & 5.34 & 11.54 & 0.1297 \\
\bottomrule\end{tabular}\end{table}

\begin{table}[t]\centering\small
\caption{B-OS-CF one-step correction $\widetilde\beta_{\rm CF}-\widehat\beta_{\rm ACV}$ and its \emph{decomposition} (CC, Scenario~2, $R=500$). The signed mean is near zero (cancellation), but the per-replicate \emph{absolute} correction is non-trivial (mean$|c|$ is $0.28$--$0.40$ of the A-CV ESD). Decomposing $\mathrm{corr}=c_{\rm split}+c_{\rm os}$ into the sample-splitting part $c_{\rm split}=\sum_k\tfrac{n_k}{n}\widehat\beta^{(-k)}_{\rm init}-\widehat\beta_{\rm ACV}$ and the pure one-step part $c_{\rm os}=\sum_k\tfrac{n_k}{n}\widehat S_k^{-1}\widehat U_k$: the orthogonal part \textbf{dominates} (mean$|c_{\rm os}|/$ESD $\approx0.28$--$0.40$ vs mean$|c_{\rm split}|/$ESD $\approx0.07$--$0.13$), so the adjustment is genuinely the orthogonal correction, not cross-fitting/splitting noise. Whether it reduces \emph{bias} is separate, addressed in the tail regime of Table~\ref{tab:sievetail}.}
\label{tab:corr}\begin{tabular}{rr r rr rr}\toprule
$n$&$\tau$&mean$|c|$&$\tfrac{|c|}{\rm ESD}$&$\tfrac{|c_{\rm os}|}{\rm ESD}$&$\tfrac{|c_{\rm split}|}{\rm ESD}$&signed mean\\\midrule
500 & 0.25 & 0.052 & 0.37 & 0.37 & 0.12 & 0.0004 \\
500 & 0.50 & 0.051 & 0.37 & 0.38 & 0.11 & 0.0033 \\
500 & 0.75 & 0.056 & 0.40 & 0.40 & 0.13 & 0.0009 \\
\addlinespace
1000 & 0.25 & 0.033 & 0.34 & 0.34 & 0.10 & -0.0022 \\
1000 & 0.50 & 0.032 & 0.35 & 0.35 & 0.09 & -0.0004 \\
1000 & 0.75 & 0.035 & 0.37 & 0.37 & 0.10 & 0.0020 \\
\addlinespace
2000 & 0.25 & 0.022 & 0.32 & 0.32 & 0.09 & -0.0011 \\
2000 & 0.50 & 0.019 & 0.28 & 0.28 & 0.07 & 0.0011 \\
2000 & 0.75 & 0.020 & 0.28 & 0.28 & 0.08 & 0.0010 \\
\bottomrule\end{tabular}\end{table}


\subsection{Surface recovery and density-weighted centering}
\begin{table}[t]\centering\small
\caption{Surface-heterogeneity DGP $\sigma(u,t,X)=\sigma_0+\sigma_1 X_1+\sigma_2(t-u)+\sigma_3(u/t)$, so $m_{\tau,0}(u,t)$ is \emph{non-parallel} across $\tau$ (\textbf{fitted estimator: CC--A-CV}; $R=500$, coverage over $300$, $B=300$). ``het'' is the base design; ``stress'' adds high truncation, $32\%$ censoring, $\sim1.3$ visits/subject. \textbf{Both} slopes are reported: bias and bootstrap coverage of $\beta_{\tau,1}$ \emph{and} the continuous-covariate $\beta_{\tau,2}$. IMSE columns are Monte-Carlo \emph{means}: IMSE$_{\rm unm}$ over the whole grid, and---now shown side by side---IMSE$_{\rm msk_v}$ over the \emph{visit}-estimable cells and IMSE$_{\rm msk_s}$ over the \emph{distinct-subject}-estimable cells. The two local-support masks: ``\%n.e.\ visit'' (fewer than $10$ pre-onset \emph{visits} within $\pm0.5$ in $(u,t)$) and the surface-relevant ``\%n.e.\ subj'' (fewer than $5$ \emph{distinct subjects}); both are replicate-specific and design-only (not strictly comparable across $n$, and the two thresholds are not information-equivalent---the subject mask is looser by construction). Because visits cluster within subject, the subject mask flags about \emph{half} as many cells as the visit mask (stress: $\approx21\%$ vs.\ $\approx41\%$), yet the two masked IMSEs are nearly identical (columns msk$_{\rm v}$ vs.\ msk$_{\rm s}$, both $\approx0.10$ under stress), so the two operational masks give \emph{qualitatively similar} conclusions---not a formal mask-invariance claim, since the thresholds are not information-matched (the subject mask is looser) and each mask is replicate-specific (a fixed common-domain IMSE would be needed to compare convergence across $n$). ESS$^{\rm ip}/n_T$ is a \emph{design diagnostic only} (the corresponding IPCW weights; the fitted CC estimator has $w\equiv1$). $\beta$ stays broadly covered (finite-sample bias rises somewhat in the smallest stress sample); the low-information boundary surface should be read only on the estimable domain.}
\label{tab:surface}\begin{tabular}{lrr rr rr rrr rr r}\toprule
&&&\multicolumn{2}{c}{$\beta_1$}&\multicolumn{2}{c}{$\beta_2$}&\multicolumn{3}{c}{IMSE}&\multicolumn{2}{c}{\%n.e.}&\\
\cmidrule(lr){4-5}\cmidrule(lr){6-7}\cmidrule(lr){8-10}\cmidrule(lr){11-12}
set&$n$&$\tau$&bias&cov&bias&cov&unm&msk$_{\rm v}$&msk$_{\rm s}$&visit&subj&$\tfrac{\rm ESS^{ip}}{n_T}$\\\midrule
het & 500 & 0.25 & -0.001 & 0.947 & 0.001 & 0.933 & 0.0326 & 0.0323 & 0.0324 & 0.03 & 0.02 & 0.98 \\
het & 500 & 0.50 & -0.005 & 0.963 & 0.003 & 0.957 & 0.0285 & 0.0284 & 0.0284 & 0.03 & 0.02 & 0.98 \\
het & 500 & 0.75 & 0.000 & 0.943 & -0.004 & 0.927 & 0.0320 & 0.0316 & 0.0317 & 0.03 & 0.02 & 0.97 \\
\addlinespace
het & 1000 & 0.25 & 0.005 & 0.920 & -0.002 & 0.937 & 0.0159 & 0.0158 & 0.0158 & 0.00 & 0.00 & 0.98 \\
het & 1000 & 0.50 & 0.000 & 0.957 & 0.003 & 0.933 & 0.0133 & 0.0133 & 0.0133 & 0.00 & 0.00 & 0.98 \\
het & 1000 & 0.75 & -0.004 & 0.933 & -0.003 & 0.943 & 0.0158 & 0.0158 & 0.0158 & 0.00 & 0.00 & 0.98 \\
\addlinespace
het & 2000 & 0.25 & 0.003 & 0.937 & 0.001 & 0.927 & 0.0077 & 0.0077 & 0.0077 & 0.00 & 0.00 & 0.98 \\
het & 2000 & 0.50 & -0.000 & 0.933 & 0.003 & 0.957 & 0.0067 & 0.0067 & 0.0067 & 0.00 & 0.00 & 0.98 \\
het & 2000 & 0.75 & 0.003 & 0.927 & -0.000 & 0.947 & 0.0074 & 0.0074 & 0.0074 & 0.00 & 0.00 & 0.98 \\
\addlinespace
stress & 500 & 0.25 & 0.019 & 0.947 & 0.010 & 0.947 & 0.1207 & 0.0956 & 0.1006 & 0.41 & 0.21 & 0.88 \\
stress & 500 & 0.50 & 0.002 & 0.957 & -0.004 & 0.973 & 0.0976 & 0.0811 & 0.0840 & 0.40 & 0.21 & 0.86 \\
stress & 500 & 0.75 & -0.026 & 0.947 & -0.007 & 0.957 & 0.1370 & 0.0900 & 0.0952 & 0.40 & 0.21 & 0.88 \\
\addlinespace
stress & 1000 & 0.25 & 0.013 & 0.953 & -0.003 & 0.940 & 0.0522 & 0.0459 & 0.0467 & 0.14 & 0.07 & 0.90 \\
stress & 1000 & 0.50 & -0.003 & 0.940 & 0.001 & 0.953 & 0.0422 & 0.0384 & 0.0390 & 0.14 & 0.07 & 0.88 \\
stress & 1000 & 0.75 & -0.008 & 0.943 & -0.001 & 0.957 & 0.0508 & 0.0463 & 0.0470 & 0.14 & 0.07 & 0.88 \\
\addlinespace
stress & 2000 & 0.25 & 0.007 & 0.960 & 0.002 & 0.943 & 0.0231 & 0.0221 & 0.0224 & 0.04 & 0.02 & 0.87 \\
stress & 2000 & 0.50 & -0.000 & 0.943 & -0.000 & 0.937 & 0.0210 & 0.0200 & 0.0204 & 0.04 & 0.02 & 0.90 \\
stress & 2000 & 0.75 & -0.010 & 0.930 & -0.001 & 0.940 & 0.0253 & 0.0242 & 0.0246 & 0.05 & 0.02 & 0.91 \\
\bottomrule\end{tabular}\end{table}

\begin{table}[t]\centering\small
\caption{Density-weighted (DW) vs.\ ordinary (ORD) centering (Scenario~2; \textbf{all three $\tau$ fit on the same dataset per replicate}, $R=600$ per $n$; IPCW weight, \emph{oracle} $f$, B-OS-\emph{NCF}; only the centering $g$ differs---Sc.~1 has $g_{\rm DW}{=}g_{\rm ORD}$ and is omitted). DW covers better in \emph{every} cell ($\Delta_{\rm cov}=+0.01$ to $+0.03$; all discordances favour DW, $c{=}0$): the paired exact McNemar is significant in $\mathbf{8/9}$ cells unadjusted, $\mathbf{4/9}$ after Holm, $\mathbf{3/9}$ after Bonferroni (we report only these cellwise/multiplicity-adjusted results; a single global $p$ is omitted as hard to audit and unnecessary). \textbf{The effect is bread calibration, not a different estimator}: DW and ORD point estimates coincide to $\sim10^{-4}$, the DW interval is only $\approx6\%$ longer (length/SE ratio $\approx1.06$) and \emph{nests} the ORD interval in $\ge99\%$ of (rep,$\tau$). Its proper $95\%$ interval score (IS; lower$=$better; width $+$ miss penalty) is \emph{statistically indistinguishable} from ORD's: the per-replicate-clustered mean $\overline{\Delta\mathrm{IS}}=-0.0012$ ($95\%$ CI $[-0.0050,0.0025]$, includes $0$), and cellwise $\Delta\mathrm{IS}={}$IS$_{\rm DW}-{}$IS$_{\rm ORD}$ is within $\approx1.6$ MCSE of $0$ \emph{except} the smallest cell ($n{=}500,\tau{=}0.25$: $+0.011$, $2.4$ MCSE, where the wider DW interval is slightly \emph{worse}-scored). \textbf{So DW improves coverage through a modestly larger sandwich SE, with \emph{no detectable change} in the proper interval score} (the cellwise $\Delta$IS are small and not uniformly favourable). The clustered $\overline{\Delta\mathrm{IS}}$ averages $\Delta$IS over the three $\tau$ \emph{within} each dataset and treats the independent datasets (across $n$) as the clusters; the resulting CI targets the \emph{equally-weighted} average $\Delta$IS over the nine evaluated $(n,\tau)$ design cells (a simulation-design summary, not a single-distribution parameter). This is an oracle-$f$, non-cross-fitted diagnostic and does not by itself carry to the feasible B-OS-CF. Coverage $=$ plug-in sandwich.}
\label{tab:ordcenter}\begin{tabular}{rr rr rr rr r}\toprule
&&\multicolumn{2}{c}{coverage}&\multicolumn{2}{c}{int.\ length}&\multicolumn{2}{c}{int.\ score}&\\
\cmidrule(lr){3-4}\cmidrule(lr){5-6}\cmidrule(lr){7-8}
$n$&$\tau$&DW&ORD&DW&ORD&DW&ORD&$\Delta$IS\\\midrule
500 & 0.25 & 0.965 & 0.952 & 0.568 & 0.535 & 0.630 & 0.619 & 0.0112 \\
500 & 0.50 & 0.953 & 0.925 & 0.541 & 0.509 & 0.624 & 0.628 & -0.0049 \\
500 & 0.75 & 0.950 & 0.930 & 0.569 & 0.535 & 0.692 & 0.691 & 0.0003 \\
\addlinespace
1000 & 0.25 & 0.950 & 0.937 & 0.409 & 0.386 & 0.484 & 0.486 & -0.0020 \\
1000 & 0.50 & 0.953 & 0.943 & 0.389 & 0.366 & 0.474 & 0.473 & 0.0012 \\
1000 & 0.75 & 0.938 & 0.918 & 0.409 & 0.385 & 0.502 & 0.507 & -0.0046 \\
\addlinespace
2000 & 0.25 & 0.940 & 0.923 & 0.291 & 0.275 & 0.363 & 0.368 & -0.0053 \\
2000 & 0.50 & 0.948 & 0.928 & 0.276 & 0.261 & 0.327 & 0.332 & -0.0044 \\
2000 & 0.75 & 0.952 & 0.935 & 0.291 & 0.275 & 0.343 & 0.346 & -0.0022 \\
\bottomrule\end{tabular}\end{table}

\begin{table}[t]\centering\small
\caption{Studentization diagnostic for the headline CC--A-CV estimator (Scenario~2, where coverage dips concentrate; $R=400$, $B=250$ per replicate). For the standardized statistic $Z_r=(\widehat\beta_r-\beta_0)/\widehat{\mathrm{SE}}_r$ we report its skewness, excess kurtosis, and $2.5/97.5\%$ quantiles; the normal-CI miss rates by side (tr$<$: truth below the CI, $Z>1.96$; tr$>$: truth above, $Z<-1.96$); the cross-replicate correlation $\rho(|\widehat\beta-\beta_0|,\widehat{\mathrm{SE}})$ (large negative $\Rightarrow$ big errors paired with small SE, which would hide undercoverage); and the coverage of the \emph{normal}-SE, \emph{percentile}, and \emph{basic} bootstrap CIs. \textbf{At the lowest-coverage cell ($n{=}500$, $\tau{=}0.25$) $Z$ has near-Gaussian central moments} (skew $-0.01$, excess kurtosis $0.02$) with $\rho\approx0$, and normal ($0.93$) and percentile ($0.93$) CIs agree---so the dip is \emph{not} explained by skew/kurtosis, SE--error dependence, or the choice of CI form. The results \emph{suggest} a finite-sample upper-tail deviation ($z_{.975}{=}2.34>1.96$); the side-specific miss-rate difference (tr$<$ $0.043$ vs tr$>$ $0.028$, i.e.\ $0.015\pm0.013$) is imprecise and not clearly distinguishable from $0$, and the deviation shrinks with $n$. Scenario~1 (omitted) is cleaner. The diagnostic gives \emph{no} evidence that percentile intervals improve on the normal-SE interval (the basic interval is worse); both show a similar mild finite-sample shortfall, \emph{consistent with} a finite-sample upper-tail deviation of the studentized statistic. The $z_{.975}{=}2.34$ rests on only $\sim10$ tail replicates ($R{=}400$), so read it qualitatively; this is a supporting re-run (cov $\approx0.93$), not a per-replicate decomposition of the headline cell.}
\label{tab:studentized}\begin{tabular}{rr rr r@{,}l rr r rrr}\toprule
&&\multicolumn{2}{c}{$Z$ shape}&\multicolumn{2}{c}{$Z$ quant.}&\multicolumn{2}{c}{CI miss}&&\multicolumn{3}{c}{coverage}\\
\cmidrule(lr){3-4}\cmidrule(lr){5-6}\cmidrule(lr){7-8}\cmidrule(lr){10-12}
$n$&$\tau$&skew&ex.kurt&\multicolumn{2}{c}{$z_{.025/.975}$}&tr$<$&tr$>$&$\rho_{e,\rm SE}$&norm&pct&basic\\\midrule
500 & 0.25 & -0.01 & 0.02 & -2.00 & 2.34 & 0.043 & 0.028 & -0.02 & 0.930 & 0.932 & 0.902 \\
500 & 0.50 & -0.09 & 0.18 & -1.90 & 1.94 & 0.022 & 0.022 & 0.10 & 0.955 & 0.965 & 0.938 \\
500 & 0.75 & -0.09 & 0.24 & -1.96 & 1.70 & 0.020 & 0.028 & 0.00 & 0.953 & 0.965 & 0.927 \\
\addlinespace
1000 & 0.25 & 0.14 & 0.30 & -1.97 & 2.24 & 0.040 & 0.030 & 0.04 & 0.930 & 0.935 & 0.915 \\
1000 & 0.50 & -0.05 & -0.17 & -1.95 & 1.95 & 0.025 & 0.025 & 0.08 & 0.950 & 0.950 & 0.917 \\
1000 & 0.75 & -0.04 & 0.31 & -1.97 & 2.05 & 0.028 & 0.028 & -0.05 & 0.945 & 0.943 & 0.927 \\
\addlinespace
2000 & 0.25 & -0.03 & 0.25 & -1.89 & 2.15 & 0.040 & 0.025 & 0.08 & 0.935 & 0.943 & 0.925 \\
2000 & 0.50 & -0.36 & 0.42 & -2.18 & 1.99 & 0.030 & 0.030 & -0.00 & 0.940 & 0.932 & 0.935 \\
2000 & 0.75 & 0.06 & -0.09 & -1.83 & 1.82 & 0.025 & 0.020 & 0.04 & 0.955 & 0.960 & 0.932 \\
\bottomrule\end{tabular}\end{table}


\subsection{Sieve growth, tuning, and cross-validation behavior}
\begin{table}[t]\centering\small
\caption{Sieve asymptotic check WITH inference, CC joint-sieve with \emph{deterministic} $J$ (Scenario~2; $R=400$, coverage over $250$, $B=250$; $\tau\in\{0.1,0.5,0.9\}$). $\sqrt n\,$bias is shown with its Monte-Carlo SE; $\dagger$ marks cells where $|\sqrt n\,\mathrm{bias}|>1.96\,$MCSE (a \emph{detectable} scaled bias). SER$=\overline{\mathrm{SE}}_{\rm boot}/\mathrm{ESD}$ and the mean $95\%$ interval length disentangle the coverage behaviour. \textbf{At $n{=}500$ the \emph{growing}-sieve bootstrap is degenerate} (mean SER $5$--$8$, mean interval length $3$--$5$ vs.\ an ESD-implied length $\approx0.6$, \emph{at all three $\tau$ including the median}): a few resamples produce enormous $\widehat{\mathrm{SE}}$ that dominate the \emph{mean}, so the near-$0.95$ coverage is an artifact of rare very-long intervals, not valid inference. We therefore mark these cells $\ddagger$ \textbf{not interpretable} (n.i.). The distributional audit (Table~\ref{tab:sievedist}) shows the \emph{median} bootstrap SE is sound (equal to the exact sieve's) and the blow-up (SE up to $96$) is confined to $2/250$ datasets ($0.8\%$, with substantial binomial uncertainty); that audit does \emph{not} attribute these events to the measured Gram-rank, local-support, or onset-diversity diagnostics, and the precise numerical trigger remains \emph{unidentified} (it audits $\tau{=}0.5$; the same $n{=}500$ growing-sieve regime shows similar blow-ups at $\tau{=}0.1,0.9$, but we do not claim an identical mechanism across $\tau$). The growing sieve is a diagnostic only---the \emph{recommended} procedure uses the small fixed sieve, where this does not occur. For $n\ge2000$ the SER is well-behaved, and there the low-bias/low-coverage cells (e.g.\ growing $\tau{=}0.1$, $n{=}4000$: $\sqrt n\,$bias $0.52$, SER $0.96$, cov $0.896$) cannot be explained by the estimated bias ($b_K$) alone---the scaled bias is only $\approx0.13$ ESD---and appear to reflect studentization/SE distortion as well.}
\label{tab:sieve}\begin{tabular}{lrr r r@{$\,\pm\,$}l r rr r}\toprule
sieve&$\tau$&$n$&$J$&\multicolumn{2}{c}{$\sqrt n\,$bias}&$\sqrt n\,$ESD&SER&len&cov$^{\rm boot}$\\\midrule
exact & 0.10 & 500 & 4 & 0.45$^\dagger$ & 0.18 & 3.67 & 1.00 & 0.645 & 0.924 \\
exact & 0.10 & 1000 & 4 & 0.19 & 0.19 & 3.79 & 1.00 & 0.469 & 0.936 \\
exact & 0.10 & 2000 & 4 & 0.05 & 0.20 & 3.92 & 0.97 & 0.331 & 0.936 \\
exact & 0.10 & 4000 & 4 & -0.21 & 0.19 & 3.71 & 1.01 & 0.233 & 0.948 \\
\addlinespace
exact & 0.50 & 500 & 4 & -0.27 & 0.15 & 2.96 & 1.06 & 0.549 & 0.936 \\
exact & 0.50 & 1000 & 4 & 0.01 & 0.15 & 2.90 & 1.06 & 0.380 & 0.956 \\
exact & 0.50 & 2000 & 4 & -0.05 & 0.15 & 2.94 & 1.03 & 0.266 & 0.960 \\
exact & 0.50 & 4000 & 4 & 0.17 & 0.15 & 3.08 & 0.99 & 0.189 & 0.940 \\
\addlinespace
exact & 0.90 & 500 & 4 & -0.26 & 0.18 & 3.67 & 1.02 & 0.658 & 0.940 \\
exact & 0.90 & 1000 & 4 & -0.31 & 0.18 & 3.67 & 1.01 & 0.460 & 0.932 \\
exact & 0.90 & 2000 & 4 & 0.05 & 0.20 & 3.94 & 0.95 & 0.328 & 0.944 \\
exact & 0.90 & 4000 & 4 & -0.33 & 0.18 & 3.65 & 1.03 & 0.234 & 0.944 \\
\addlinespace
grow & 0.10 & 500 & 7 & 0.92$^\dagger$ & 0.18 & 3.64 & 5.22 & 3.328 & n.i.$^\ddagger$ \\
grow & 0.10 & 1000 & 7 & 0.78$^\dagger$ & 0.18 & 3.62 & 1.03 & 0.462 & 0.952 \\
grow & 0.10 & 2000 & 8 & 1.03$^\dagger$ & 0.20 & 3.93 & 0.96 & 0.331 & 0.944 \\
grow & 0.10 & 4000 & 9 & 0.52$^\dagger$ & 0.20 & 3.94 & 0.96 & 0.233 & 0.896 \\
\addlinespace
grow & 0.50 & 500 & 7 & -0.10 & 0.15 & 3.07 & 6.66 & 3.585 & n.i.$^\ddagger$ \\
grow & 0.50 & 1000 & 7 & 0.08 & 0.15 & 2.92 & 1.06 & 0.382 & 0.960 \\
grow & 0.50 & 2000 & 8 & -0.07 & 0.15 & 3.08 & 1.00 & 0.269 & 0.944 \\
grow & 0.50 & 4000 & 9 & -0.01 & 0.16 & 3.22 & 0.95 & 0.190 & 0.940 \\
\addlinespace
grow & 0.90 & 500 & 7 & -0.70$^\dagger$ & 0.18 & 3.67 & 7.64 & 4.911 & n.i.$^\ddagger$ \\
grow & 0.90 & 1000 & 7 & -0.61$^\dagger$ & 0.20 & 3.94 & 0.96 & 0.467 & 0.964 \\
grow & 0.90 & 2000 & 8 & -0.48$^\dagger$ & 0.18 & 3.54 & 1.05 & 0.325 & 0.968 \\
grow & 0.90 & 4000 & 9 & -0.46$^\dagger$ & 0.19 & 3.81 & 0.98 & 0.232 & 0.928 \\
\bottomrule\end{tabular}\end{table}

\begin{table}[t]\centering\small
\caption{Does the one-step reduce the growing-sieve tail \emph{bias}? Deterministic growing sieve, same replicates through A-joint-sieve, cross-fitted \textbf{B-OS-CF}, and a richer-sieve A-Rich ($J{+}2$) sensitivity (CC, Sc.\ 2; $R=400$, coverage over $150$, $B=120$---a \emph{pilot-grade} inference subset, coverage MCSE $\approx0.018$; $\sqrt n\,$bias MCSE $\approx0.15$--$0.20$). At \emph{both} tails B-OS-CF consistently moves the scaled bias toward zero, increasingly with $n$ (e.g.\ $\tau{=}0.1$: $0.66{\to}0.28$, $0.53{\to}0.19$; $\tau{=}0.9$: $-0.49{\to}-0.12$ at $n{=}4000$), whereas the fixed-$(J{+}2)$ richer sieve \emph{worsens} it; paired Monte-Carlo intervals would be needed to resolve which individual moves are significant. \textbf{However}, SER and RMSE show a \emph{bias}-targeted correction only: B-OS-CF's SER $\approx0.93$--$0.97$, coverage is not improved over the deterministic A-joint-sieve (\emph{not} the CV-selected A-CV; the comparator here is the fixed growing sieve), and RMSE is essentially unchanged (the bias drop is offset by a small variance rise)---the scaled bias is already small relative to the SE here, so debiasing is a safeguard, not a coverage fix at these $n$.}
\label{tab:sievetail}\begin{tabular}{rr r rrrr rrrr rr}\toprule
&&& \multicolumn{4}{c}{A-joint-sieve} & \multicolumn{4}{c}{B-OS-CF} & \multicolumn{2}{c}{A-Rich}\\
\cmidrule(lr){4-7}\cmidrule(lr){8-11}\cmidrule(lr){12-13}
$\tau$&$n$&$J$&$\sqrt n\,$bias&SER&RMSE&cov&$\sqrt n\,$bias&SER&RMSE&cov&$\sqrt n\,$bias&cov\\\midrule
0.10 & 1000 & 7 & 1.09 & 0.98 & 0.126 & 0.953 & 1.01 & 0.93 & 0.128 & 0.913 & 1.36 & 0.907 \\
0.10 & 2000 & 8 & 0.66 & 0.96 & 0.088 & 0.973 & 0.28 & 0.93 & 0.089 & 0.980 & 0.91 & 0.960 \\
0.10 & 4000 & 9 & 0.53 & 0.97 & 0.062 & 0.947 & 0.19 & 0.96 & 0.064 & 0.940 & 0.68 & 0.920 \\
\addlinespace
0.50 & 1000 & 7 & 0.17 & 1.03 & 0.095 & 0.947 & 0.16 & 1.02 & 0.095 & 0.940 & 0.16 & 0.953 \\
0.50 & 2000 & 8 & -0.09 & 0.97 & 0.070 & 0.940 & -0.15 & 0.96 & 0.072 & 0.953 & -0.12 & 0.933 \\
0.50 & 4000 & 9 & -0.13 & 1.00 & 0.048 & 0.940 & -0.08 & 0.99 & 0.050 & 0.953 & -0.14 & 0.947 \\
\addlinespace
0.90 & 1000 & 7 & -0.88 & 0.98 & 0.125 & 0.967 & -0.75 & 0.95 & 0.126 & 0.973 & -1.23 & 0.967 \\
0.90 & 2000 & 8 & -0.16 & 1.02 & 0.083 & 0.947 & 0.09 & 0.95 & 0.089 & 0.927 & -0.33 & 0.940 \\
0.90 & 4000 & 9 & -0.49 & 1.02 & 0.059 & 0.967 & -0.12 & 0.97 & 0.062 & 0.960 & -0.67 & 0.967 \\
\bottomrule\end{tabular}\end{table}

\begin{table}[t]\centering\small
\caption{Tuning sensitivity to the sparsity bandwidth $a_n\in\{.03,.05,.08\}$ and projection dimension $df_g\in\{3,4,5\}$ (defaults $.05,4$), now over the FULL grid $n\in\{500,1000,2000\}\times\tau\in\{.1,.5,.9\}$ (Sc.\ 2, $R=400$). Per $(n,\tau)$ we report the \emph{range} across the $9$ combinations of: B-OS-CF $\beta_{\tau,1}$ bias and ESD; the plug-in sandwich SER and its coverage; the worst $\widehat f$-cap rate; and the median bread condition number. \textbf{Point estimation is robust to tuning everywhere} (bias range $\le0.03$, ESD range $\le0.03$, even at $n{=}500$ and the tails). \textbf{Inference is not}: at $n{=}500$ the sandwich SER ranges $0.57$--$1.03$ (coverage $0.69$--$0.93$) across the grid---confirming that the $a_n$-dependent $\widehat f$ destabilizes the analytic SE at small $n$ (the bootstrap is used for inference instead). This tuning-sensitivity of the analytic sandwich \emph{persists at the tails even at} $n{=}2000$ (for $\tau{=}0.1,0.9$, SER ranges $0.89$--$1.24$, coverage $0.91$--$0.98$ across the grid); only the central quantile is comparatively stable there. Note this varies the \emph{sandwich}, whereas the recommended procedure uses the bootstrap---so it bounds the analytic SE's tuning-robustness, not the bootstrap's. $\widehat f$-cap rates are $\le2.2\%$ and $\widehat S$ is well-conditioned ($\kappa\le7$) throughout.}
\label{tab:tuning}\begin{tabular}{rr rr rr rr}\toprule
&&\multicolumn{2}{c}{bias / ESD range}&\multicolumn{2}{c}{SER / cov range}&&\\
\cmidrule(lr){3-4}\cmidrule(lr){5-6}
$n$&$\tau$&bias rng&ESD rng&SER rng&cov rng&$\widehat f$-cap&$\kappa_{\max}$\\\midrule
500 & 0.10 & [-0.006,0.015] & [0.176,0.203] & [0.57,1.02] & [0.71,0.93] & 0.022 & 7 \\
500 & 0.50 & [-0.015,0.013] & [0.143,0.153] & [0.68,0.95] & [0.76,0.93] & 0.021 & 6 \\
500 & 0.90 & [-0.019,0.002] & [0.177,0.207] & [0.61,1.03] & [0.69,0.92] & 0.022 & 6 \\
1000 & 0.10 & [-0.009,0.018] & [0.130,0.142] & [0.71,1.17] & [0.79,0.97] & 0.009 & 5 \\
1000 & 0.50 & [-0.009,0.012] & [0.096,0.106] & [0.84,1.03] & [0.89,0.95] & 0.008 & 5 \\
1000 & 0.90 & [-0.010,0.006] & [0.128,0.144] & [0.73,1.19] & [0.80,0.98] & 0.008 & 5 \\
2000 & 0.10 & [-0.010,0.009] & [0.088,0.099] & [0.91,1.24] & [0.93,0.98] & 0.003 & 5 \\
2000 & 0.50 & [-0.005,0.002] & [0.069,0.074] & [0.90,1.01] & [0.92,0.95] & 0.003 & 4 \\
2000 & 0.90 & [-0.009,0.002] & [0.086,0.104] & [0.89,1.21] & [0.91,0.98] & 0.003 & 4 \\
\bottomrule\end{tabular}\end{table}

\begin{table}[t]\centering\small
\caption{Variance-ladder for the plug-in clustered sandwich SER ($=\overline{\mathrm{SE}}/$Monte-Carlo ESD; $R=500$). $\widehat f/f_0$ is the per-cell \emph{median over replicates and rows} of the sparsity ratio. Rungs replace plug-in nuisances by truth one at a time, holding the A-CV point estimator fixed; the $g$-true rung uses the \emph{finite-sieve} oracle $g_J$ at the point-estimator's df; fix-$J$ is a \emph{separate} rung that re-selects the point estimator at a deterministic df (normalized by its own ESD). Both the $f$-true ($0.88$--$1.07$) and oracle-$(f,g)$ ($0.89$--$1.06$) rungs are near nominal, while the $g$-true rung \emph{alone} stays low ($0.54$--$0.95$): replacing the sparsity $\widehat f$ nearly closes the gap, replacing $g$ alone does not. Sparsity estimation is thus the \emph{largest identified contributor} in the small- and moderate-$n$ cells, though residual distortion remains and the $f$-true improvement is not monotone in every large-$n$ cell; the bread nuisances must be replaced \emph{jointly} to be consistent. Row/cluster $<1$ confirms clustering is already included (and matters).}
\label{tab:ladder}\begin{tabular}{rrr r rrrrr r}\toprule
Sc.&$n$&$\tau$&$\widehat f/f_0$&plug&$f$-true&$g$-true&oracle&fix-$J$&$\tfrac{\rm row}{\rm clus}$\\\midrule
1 & 500 & 0.25 & 1.04 & 0.70 & 0.88 & 0.54 & 0.89 & 0.59 & 0.63 \\
1 & 500 & 0.50 & 1.04 & 0.77 & 0.96 & 0.60 & 0.97 & 0.67 & 0.61 \\
1 & 500 & 0.75 & 1.02 & 0.75 & 0.92 & 0.58 & 0.93 & 0.62 & 0.63 \\
\addlinespace
1 & 1000 & 0.25 & 1.02 & 0.85 & 0.95 & 0.74 & 0.95 & 0.79 & 0.62 \\
1 & 1000 & 0.50 & 1.02 & 0.89 & 0.98 & 0.79 & 0.98 & 0.82 & 0.60 \\
1 & 1000 & 0.75 & 1.02 & 0.87 & 0.96 & 0.76 & 0.96 & 0.81 & 0.62 \\
\addlinespace
1 & 2000 & 0.25 & 1.00 & 0.93 & 0.97 & 0.88 & 0.97 & 0.90 & 0.60 \\
1 & 2000 & 0.50 & 1.00 & 1.02 & 1.07 & 0.95 & 1.06 & 0.99 & 0.59 \\
1 & 2000 & 0.75 & 1.00 & 0.94 & 0.98 & 0.87 & 0.98 & 0.91 & 0.60 \\
\addlinespace
2 & 500 & 0.25 & 1.03 & 0.78 & 0.91 & 0.55 & 0.92 & 0.68 & 0.63 \\
2 & 500 & 0.50 & 1.02 & 0.80 & 0.93 & 0.64 & 0.95 & 0.69 & 0.61 \\
2 & 500 & 0.75 & 1.04 & 0.83 & 0.98 & 0.64 & 0.99 & 0.72 & 0.64 \\
\addlinespace
2 & 1000 & 0.25 & 1.01 & 0.93 & 0.99 & 0.82 & 1.00 & 0.87 & 0.61 \\
2 & 1000 & 0.50 & 1.02 & 0.86 & 0.93 & 0.78 & 0.93 & 0.82 & 0.60 \\
2 & 1000 & 0.75 & 1.01 & 0.89 & 0.95 & 0.79 & 0.95 & 0.84 & 0.61 \\
\addlinespace
2 & 2000 & 0.25 & 1.00 & 0.98 & 1.00 & 0.92 & 1.00 & 0.95 & 0.60 \\
2 & 2000 & 0.50 & 1.01 & 0.98 & 1.01 & 0.91 & 1.01 & 0.96 & 0.59 \\
2 & 2000 & 0.75 & 1.00 & 0.95 & 0.97 & 0.88 & 0.96 & 0.92 & 0.60 \\
\bottomrule\end{tabular}\end{table}

\begin{table}[t]\centering\small
\caption{Frequency of the CV-selected spline complexity $\widehat J\in\{3,4,5\}$ (the A-CV tuning), by Scenario, $n$, and $\tau$ ($R=500$ per cell). Here $\widehat J$ is the per-margin \texttt{df}~$=$~degree~$+$~(interior knots), so $\widehat J{=}3$ is a genuine \emph{cubic} with \emph{zero} interior knots ($4$ B-spline functions/margin, $16$ tensor functions), not a low-degree fit; the adaptive choice is the number of interior knots (defined in the reproducibility specification). \textbf{The selection concentrates on the grid minimum $\widehat J=3$} ($81$--$86\%$), with $P(\widehat J{=}4)\approx0.10$--$0.15$ and $P(\widehat J{=}5)\approx0.02$--$0.04$, and is nearly invariant in $n$ and $\tau$. Two consequences: (i) A-CV is here close to a \emph{fixed} $J{=}3$ estimator, so its ``adaptive'' component is small and the A-Rich ($J{+}2$) contrast is effectively $J{=}5$ vs.\ $J{=}3$ rather than a perturbation around an interior optimum; (ii) the selection saturates at the \emph{minimum} complexity available within the cubic-spline family ($\widehat J{=}3$ is zero interior knots, which a wider $J$ grid cannot undercut), so whether a lower-degree or parametric surface would be preferred is not determined by this experiment. This tempers the adaptive-tuning rung of Table~\ref{tab:ladder}.}
\label{tab:cvfreq}\begin{tabular}{rrr rrr r}\toprule
Sc.&$n$&$\tau$&$P(\widehat J{=}3)$&$P(\widehat J{=}4)$&$P(\widehat J{=}5)$&mode\\\midrule
1 & 500 & 0.25 & 0.81 & 0.15 & 0.04 & 3 \\
1 & 500 & 0.50 & 0.84 & 0.13 & 0.03 & 3 \\
1 & 500 & 0.75 & 0.86 & 0.12 & 0.03 & 3 \\
\addlinespace
1 & 1000 & 0.25 & 0.85 & 0.10 & 0.04 & 3 \\
1 & 1000 & 0.50 & 0.83 & 0.13 & 0.04 & 3 \\
1 & 1000 & 0.75 & 0.82 & 0.14 & 0.04 & 3 \\
\addlinespace
1 & 2000 & 0.25 & 0.83 & 0.14 & 0.04 & 3 \\
1 & 2000 & 0.50 & 0.85 & 0.12 & 0.03 & 3 \\
1 & 2000 & 0.75 & 0.82 & 0.15 & 0.03 & 3 \\
\addlinespace
2 & 500 & 0.25 & 0.82 & 0.14 & 0.04 & 3 \\
2 & 500 & 0.50 & 0.85 & 0.13 & 0.02 & 3 \\
2 & 500 & 0.75 & 0.82 & 0.14 & 0.04 & 3 \\
\addlinespace
2 & 1000 & 0.25 & 0.83 & 0.14 & 0.04 & 3 \\
2 & 1000 & 0.50 & 0.84 & 0.12 & 0.04 & 3 \\
2 & 1000 & 0.75 & 0.85 & 0.13 & 0.02 & 3 \\
\addlinespace
2 & 2000 & 0.25 & 0.84 & 0.11 & 0.04 & 3 \\
2 & 2000 & 0.50 & 0.84 & 0.13 & 0.03 & 3 \\
2 & 2000 & 0.75 & 0.86 & 0.11 & 0.03 & 3 \\
\bottomrule\end{tabular}\end{table}

\begin{table}[t]\centering\small
\caption{Is the data-adaptive A-CV materially different from a \emph{fixed} small sieve? (Scenario~2, CC estimator, same replicates; $R=400$, $B=200$.) Since CV selects $\widehat J{=}3$ in $80$--$86\%$ of replicates (Table~\ref{tab:cvfreq}), we put A-CV (CV-selected $\widehat J$) head-to-head with fixed $J{=}3$. \textbf{CV genuinely prefers $J{=}3$---it is not tie-breaking noise}: the normalized CV-loss gaps (gap$_J=(\mathrm{CV}(J){-}\mathrm{CV}(3))/|\mathrm{CV}(3)|$, the relative increase in held-out \emph{complete-event} (CC, $\omega{=}1$) check loss---the CC member tunes its df by CC-CV, so \emph{no} $\widehat G_C$ enters the CC pipeline) are positive (gap$_4\approx0.004$--$0.029$, $P(\text{gap}_4{>}0)\approx0.85$), shrinking with $n$. \textbf{And A-CV $\approx$ fixed $J{=}3$}: the point estimates differ by only $\overline{|\Delta|}\approx0.002$--$0.007$, with near-identical bias, ESD, and bootstrap coverage. On the \emph{surface} IMSE, however, fixed $J{=}3$ is \emph{consistently lower}---by $\approx9$--$13\%$ in the displayed cells; the paired $\overline{\Delta\mathrm{IMSE}}={}$IMSE$_{\rm ACV}-{}$IMSE$_{J3}$ (last column, MCSE in parentheses) is positive throughout, consistent with occasional overfitting when A-CV selects $J{\in}\{4,5\}$. So in these DGPs \textbf{CV mostly reproduces the fixed $J{=}3$ fit and offers little measurable gain on $\beta$}---the headline simulation therefore essentially validates a small fixed sieve. We do not claim CV equals the per-replicate oracle complexity, nor that it ``protects'' against richer fits (the average CV criterion and the majority of replicates favour the smallest candidate, but $\approx14$--$20\%$ select $J{\in}\{4,5\}$); only that its preference for $J{=}3$ is systematic, not tie-breaking noise.}
\label{tab:fixedj}\begin{tabular}{rr r rr r rr rr r}\toprule
&&&\multicolumn{2}{c}{CV gap}&&\multicolumn{2}{c}{coverage}&\multicolumn{2}{c}{IMSE}&\\
\cmidrule(lr){4-5}\cmidrule(lr){7-8}\cmidrule(lr){9-10}
$n$&$\tau$&$P(\widehat J{=}3)$&$\substack{\rm CV4\\-\rm CV3}$&$\substack{\rm CV5\\-\rm CV3}$&$\overline{|\Delta\beta|}$&A-CV&$J{=}3$&A-CV&$J{=}3$&$\Delta$IMSE\\\midrule
500 & 0.25 & 0.81 & 0.026 & 0.074 & 0.0064 & 0.927 & 0.935 & 0.0335 & 0.0295 & 0.0040(0.0006) \\
500 & 0.50 & 0.86 & 0.023 & 0.059 & 0.0033 & 0.950 & 0.953 & 0.0273 & 0.0249 & 0.0024(0.0004) \\
500 & 0.75 & 0.82 & 0.026 & 0.065 & 0.0062 & 0.950 & 0.955 & 0.0334 & 0.0300 & 0.0035(0.0006) \\
\addlinespace
1000 & 0.25 & 0.82 & 0.011 & 0.027 & 0.0033 & 0.930 & 0.932 & 0.0160 & 0.0142 & 0.0019(0.0003) \\
1000 & 0.50 & 0.81 & 0.009 & 0.020 & 0.0027 & 0.945 & 0.943 & 0.0138 & 0.0120 & 0.0018(0.0002) \\
1000 & 0.75 & 0.81 & 0.012 & 0.025 & 0.0032 & 0.943 & 0.943 & 0.0153 & 0.0138 & 0.0015(0.0002) \\
\addlinespace
2000 & 0.25 & 0.83 & 0.005 & 0.010 & 0.0018 & 0.938 & 0.935 & 0.0078 & 0.0070 & 0.0008(0.0001) \\
2000 & 0.50 & 0.83 & 0.004 & 0.009 & 0.0013 & 0.940 & 0.943 & 0.0063 & 0.0056 & 0.0007(0.0001) \\
2000 & 0.75 & 0.83 & 0.005 & 0.011 & 0.0018 & 0.955 & 0.958 & 0.0077 & 0.0069 & 0.0008(0.0001) \\
\bottomrule\end{tabular}\end{table}


\subsection{Fold discipline and randomization}
\begin{table}[t]\centering\small
\caption{Fold-randomization sensitivity of the CC \textbf{B-OS-CF} bootstrap (Scenario~2; $R=250$, $B=120$, $M=8$ split draws). The point estimator is the single \emph{realized} fold split (as in the headline) or the \emph{average} over $M$ splits; the bootstrap SE either redraws folds each resample (headline default) or keeps each original subject's realized fold (\emph{fixed}). SER is relative to the matching point estimator's ESD. \textbf{Redrawing folds inflates the SER by only $\approx0.01$--$0.015$ over fixing them} ($\Delta$SER $=$ SER$_{\rm redr}-$SER$_{\rm fix}$; paired MCSE in parentheses $=$ the sd of the per-replicate $\widehat{\mathrm{SE}}_{\rm redr}-\widehat{\mathrm{SE}}_{\rm fix}$ divided by the common ESD and $\sqrt R$, treating ESD as fixed); what matters is the \emph{effect size}---fold redrawing adds only $\approx0.01$--$0.015$ to the SER, a small contribution---so the bulk of the mild B-OS-CF conservativeness is the duplicate-ID fold-leakage correction (Table~\ref{tab:leakfold}), not fold randomization. The headline bootstrap \emph{redraws} folds each resample and so marginalizes over fold randomness---this is the recommended inference; the fixed-fold and repeated-split rows are \emph{sensitivity checks}. Fixed folds are marginally tighter and empirically well calibrated here (coverage $0.948$--$0.968$); averaging the point over splits lowers its ESD but, with the redrawn SE, \emph{over}-covers ($\le0.98$). Per-cell coverage MCSE $\approx0.014$ ($R{=}250$), so cell-level coverage gaps should not be over-read.}
\label{tab:foldsens}\begin{tabular}{rr rr rr r rrr}\toprule
&&\multicolumn{2}{c}{ESD}&\multicolumn{2}{c}{SER}&&\multicolumn{3}{c}{coverage}\\
\cmidrule(lr){3-4}\cmidrule(lr){5-6}\cmidrule(lr){8-10}
$n$&$\tau$&real&avg&fix&redr&$\Delta$SER&$\substack{\rm real\\ \rm fix}$&$\substack{\rm real\\ \rm redr}$&$\substack{\rm avg\\ \rm redr}$\\\midrule
500 & 0.25 & 0.160 & 0.147 & 1.03 & 1.04 & 0.012(0.005) & 0.944 & 0.952 & 0.960 \\
500 & 0.50 & 0.153 & 0.137 & 1.02 & 1.03 & 0.013(0.005) & 0.956 & 0.968 & 0.980 \\
500 & 0.75 & 0.157 & 0.144 & 1.05 & 1.07 & 0.014(0.005) & 0.948 & 0.944 & 0.968 \\
\addlinespace
1000 & 0.50 & 0.109 & 0.101 & 1.00 & 1.01 & 0.013(0.005) & 0.968 & 0.964 & 0.976 \\
\addlinespace
2000 & 0.50 & 0.073 & 0.068 & 1.05 & 1.07 & 0.016(0.005) & 0.948 & 0.948 & 0.952 \\
\bottomrule\end{tabular}\end{table}

\begin{table}[t]\centering\small
\caption{\textbf{Controlled implementation comparison} of the duplicate-ID fold-leakage correction (CC B-OS-CF, Scenario~2; $R=250$, $B=150$). On the \emph{same} datasets and resamples the bootstrap SE is computed two ways: \emph{leaked} (the previous implementation---fold by the resampled id, so duplicate copies of a subject split across folds and one trains a nuisance the other is scored against) vs.\ \emph{clean} (fold by orig\_id, duplicates kept together). The last column is the \emph{paired} ratio $\widehat{\mathrm{SE}}_{\rm clean}/\widehat{\mathrm{SE}}_{\rm leaked}$ computed per replicate then averaged (MCSE in parentheses $=$ sd of the per-replicate ratio $/\sqrt R$). \textbf{The leak systematically tightens the bootstrap SE by $6$--$11\%$}: SER$_{\rm leaked}<1$ in nearly every cell with coverage dipping to $\approx0.93$, whereas enforcing fold discipline restores SER$_{\rm clean}\ge1$ and near-nominal coverage ($0.94$--$0.98$). This attributes most of the mild B-OS-CF conservativeness of Table~\ref{tab:estsim} to the correction; combined with the $\le0.015$ fold-\emph{randomization} contribution of Table~\ref{tab:foldsens}, it explains the bulk of the shift from underestimation to mild conservativeness, though residual overcoverage in several cells (clean SER up to $1.17$) is \emph{not} uniquely attributable to it. Per-cell coverage MCSE $\approx0.014$ ($R{=}250$), so cell-level coverage gaps ($\approx0.016$--$0.020$) should not be over-read individually.}
\label{tab:leakfold}\begin{tabular}{rr rr r rr}\toprule
&&\multicolumn{2}{c}{SER}&&\multicolumn{2}{c}{coverage}\\
\cmidrule(lr){3-4}\cmidrule(lr){6-7}
$n$&$\tau$&leaked&clean&$\tfrac{\rm clean}{\rm leaked}$&leaked&clean\\\midrule
500 & 0.25 & 0.974 & 1.044 & 1.076(0.004) & 0.932 & 0.952 \\
500 & 0.50 & 0.976 & 1.036 & 1.063(0.004) & 0.940 & 0.968 \\
500 & 0.75 & 0.999 & 1.071 & 1.074(0.005) & 0.928 & 0.948 \\
\addlinespace
1000 & 0.25 & 0.977 & 1.075 & 1.102(0.004) & 0.956 & 0.972 \\
1000 & 0.50 & 0.931 & 1.021 & 1.099(0.004) & 0.944 & 0.972 \\
1000 & 0.75 & 1.033 & 1.133 & 1.098(0.004) & 0.936 & 0.952 \\
\addlinespace
2000 & 0.25 & 0.956 & 1.057 & 1.108(0.005) & 0.932 & 0.948 \\
2000 & 0.50 & 0.978 & 1.068 & 1.093(0.004) & 0.932 & 0.940 \\
2000 & 0.75 & 1.065 & 1.171 & 1.100(0.004) & 0.976 & 0.984 \\
\bottomrule\end{tabular}\end{table}


\subsection{Non-Gaussian errors and crossing}
\begin{table}[t]\centering\small
\caption{Distributional robustness on the \emph{heteroscedastic} Scenario~2 DGP ($\sigma_1=0.4$), so the truth is $\beta_{\tau,1}^{\rm true}=\beta_{01}+\sigma_1 q_\tau(F_G)$ and shifts with the error law $F_G$ (reported as $\beta_1^{\rm tr}$ for independent verification). The marker error is a standardized non-normal law via a Gaussian copula (latent-Gaussian dependence preserved; raw within-subject correlation not held fixed): ``t5'' is standardized Student-$t_5$, ``skew'' a standardized right-skew lognormal. $R=500$, coverage over $300$ ($B=300$; MC SE $\approx0.013$). Point estimation stays approximately unbiased for both slopes; bootstrap coverage is generally satisfactory ($0.91$--$0.98$), with moderate under-coverage (min $\approx0.91$) in a few cells.}
\label{tab:dist}\begin{tabular}{lrr r rr r rr}\toprule
err&$n$&$\tau$&$\beta_1^{\rm tr}$&bias$_{\beta1}$&ESD$_{\beta1}$&cov$_{\beta1}$&bias$_{\beta2}$&cov$_{\beta2}$\\\midrule
skew & 500 & 0.25 & 0.222 & 0.004 & 0.084 & 0.960 & -0.001 & 0.963 \\
skew & 500 & 0.50 & 0.412 & -0.004 & 0.110 & 0.970 & 0.002 & 0.950 \\
skew & 500 & 0.75 & 0.677 & -0.014 & 0.165 & 0.947 & -0.001 & 0.910 \\
\addlinespace
skew & 1000 & 0.25 & 0.222 & 0.002 & 0.063 & 0.913 & 0.001 & 0.957 \\
skew & 1000 & 0.50 & 0.412 & -0.001 & 0.076 & 0.980 & -0.001 & 0.930 \\
skew & 1000 & 0.75 & 0.677 & -0.004 & 0.118 & 0.957 & 0.002 & 0.943 \\
\addlinespace
skew & 2000 & 0.25 & 0.222 & 0.001 & 0.040 & 0.977 & 0.001 & 0.957 \\
skew & 2000 & 0.50 & 0.412 & -0.003 & 0.057 & 0.923 & -0.001 & 0.953 \\
skew & 2000 & 0.75 & 0.677 & -0.006 & 0.083 & 0.940 & -0.001 & 0.963 \\
\addlinespace
t5 & 500 & 0.25 & 0.275 & 0.007 & 0.131 & 0.953 & 0.001 & 0.950 \\
t5 & 500 & 0.50 & 0.500 & -0.002 & 0.112 & 0.947 & -0.002 & 0.967 \\
t5 & 500 & 0.75 & 0.725 & -0.004 & 0.118 & 0.960 & -0.000 & 0.957 \\
\addlinespace
t5 & 1000 & 0.25 & 0.275 & 0.008 & 0.090 & 0.930 & 0.001 & 0.950 \\
t5 & 1000 & 0.50 & 0.500 & 0.005 & 0.077 & 0.950 & 0.004 & 0.943 \\
t5 & 1000 & 0.75 & 0.725 & 0.005 & 0.086 & 0.960 & 0.000 & 0.960 \\
\addlinespace
t5 & 2000 & 0.25 & 0.275 & 0.005 & 0.066 & 0.947 & 0.000 & 0.957 \\
t5 & 2000 & 0.50 & 0.500 & 0.001 & 0.058 & 0.927 & 0.001 & 0.943 \\
t5 & 2000 & 0.75 & 0.725 & -0.003 & 0.062 & 0.940 & 0.001 & 0.937 \\
\bottomrule\end{tabular}\end{table}

\begin{table}[t]\centering\small
\caption{Distributional audit of the deterministic-sieve subject bootstrap at $\tau{=}0.5$ (same DGP/seed as Table~\ref{tab:sieve}; $R=250$, $B=150$), with a \emph{scale-invariant} numerical-stability diagnostic. \textbf{The median bootstrap SE equals the exact sieve's at every $n$}; only the $n{=}500$ growing sieve blows up---SE$_{\max}=96$, and the SE$>3\times$med fraction is $0.8\%$ ($2/250$ datasets). To test \emph{why}, we replace the raw condition number (which is huge even where the SE is stable) with scale-invariant degeneracy of the standardized design Gram of $[X_1,X_2,\text{surface basis}]$: the \emph{effective rank} $r_{\rm eff}=(\sum_j\lambda_j)^2/\sum_j\lambda_j^2$ (a participation ratio over the Gram eigenvalues, \emph{not} a numerical rank). Near-zero-support surface columns are removed by a deterministic pre-fit screen; \emph{all} columns surviving that screen are kept in the QR fit (no further numerical-rank truncation, and QR never fails). $N_{\rm low}=$ the count of surface columns carrying $<3$ rows of mass. The $r_{\rm eff}$ denominator is the design width $=p{+}K_{\rm ret}$, the $p{=}2$ covariates plus the retained surface tensor columns (near-zero columns dropped): $27{=}2{+}25$ at $J{=}4$, $63{=}2{+}61$ of $64$ at $J{=}7$. The basis \emph{is} over-rich at $n{=}500$ ($r_{\rm eff}\approx21$ of $63$ columns, $\approx8$ low-support columns)---\textbf{but this does \emph{not} explain the blow-up}: $\mathrm{corr}(\log\widehat{\mathrm{SE}},-\log\lambda_{\min})\approx0$, and the $2$ extreme-SE datasets are \emph{not} among the low-$\lambda_{\min}$ (weakest-Gram) ones (coincidence $0/2$). We also audited the blow-up at the bootstrap-\emph{resample} level: the extreme-$|\widehat\beta^*|$ resamples (top $0.5\%$) are \emph{indistinguishable} from the bulk in lost-support columns ($7.7$ vs $7.6$), standardized-Gram $\lambda_{\min}$ ($1.5$ vs $1.6{\times}10^{-3}$), and onset diversity ($191$ vs $195$), with all correlations $\approx0$. So the instability is \emph{not} a Gram-rank or support failure at either the dataset or the resample level; it appears to be a rare numerical/optimization instability of the over-rich small-$n$ sieve rather than a failure of the measured design-rank diagnostics, and we leave the precise trigger \emph{unidentified} (it does not affect the recommended fixed sieve). We therefore present it descriptively (a small-$n$ over-rich-sieve artifact, $2/250$ datasets) and \emph{withdraw} any condition-number causal attribution. The recommended fixed sieve showed no analogous instability in the audited runs. (Column big/$\cap$rd: the number of datasets with bootstrap SE $>3\times$ the median, and how many of those also fall in the weakest-Gram decile ($\lambda_{\min}$ in the bottom $10\%$; ``weakest-Gram'', not necessarily rank-deficient).)}
\label{tab:sievedist}\begin{tabular}{l r r r r r r r r}\toprule
&&&\multicolumn{2}{c}{bootstrap SE}&\multicolumn{2}{c}{degeneracy}&&\\
\cmidrule(lr){4-5}\cmidrule(lr){6-7}
sieve&$n$&$J$&med&max&$r_{\rm eff}$&$N_{\rm low}$&$\rho_{\log{\rm SE},-\log\lambda}$&big/$\cap$rd\\\midrule
exact & 500 & 4 & 0.138 & 0.20 & 7/27 & 2 & 0.03 & 0/0 \\
exact & 1000 & 4 & 0.097 & 0.13 & 7/27 & 2 & 0.04 & 0/0 \\
exact & 2000 & 4 & 0.067 & 0.09 & 7/27 & 2 & 0.02 & 0/0 \\
\addlinespace
grow & 500 & 7 & 0.139 & 96.44 & 21/63 & 8 & -0.03 & 2/0 \\
grow & 1000 & 7 & 0.097 & 0.12 & 20/63 & 7 & -0.02 & 0/0 \\
grow & 2000 & 8 & 0.069 & 0.08 & 26/77 & 8 & 0.06 & 0/0 \\
\bottomrule\end{tabular}\end{table}

\begin{table}[t]\centering\small
\caption{Quantile-crossing on a \emph{dense} $30\times30$ interior $(u,t)$ grid $\times$ $10$ covariate profiles ($\tau\in\{.1,.25,.5,.75,.9\}$, adjacent pairs; $R=300$). Crossing rates split interior/boundary; ``$\bar u/t$'' is the mean $u/t$ of crossing \emph{incidents} (small $\Rightarrow$ near the low-age boundary). The interior crossing rate is $0$ in every cell; boundary crossing rates round to $\approx0$ by $n=1000$, but at $n{=}500$ a non-trivial \emph{fraction of replicates} ($\le15\%$, Sc.~2) show \emph{some} boundary crossing with per-incident magnitude up to $0.65$---so we rearrange by default. cal$_{\rm in}$ is the in-sample $\tau$-calibration error; cal$^{\rm OOS}_{\rm vis}$/cal$^{\rm OOS}_{\rm sub}$ are the \emph{out-of-sample} errors on an \emph{independent test draw} (same DGP), \emph{visit}-weighted (mean over $\approx620$ test rows) vs.\ \emph{subject}-equal-weighted (so multi-visit subjects do not dominate). The two agree ($0.011$--$0.022$); cal$^{\rm OOS}_{\rm vis}$ is the value \emph{before} rearrangement, and rearrangement shifts it by at most $0.001$ in every cell (largest, Sc.~2 $n{=}500$: $0.0194\!\to\!0.0184$; the full before/after pair is in the released \texttt{crossing.csv}), so rearrangement does not distort out-of-sample calibration under either weighting.}
\label{tab:crossing}\begin{tabular}{rr rr r r r r rr}\toprule
Sc.&$n$&cr$_{\rm int}$&cr$_{\rm bnd}$&frac.\ reps&max mag&$\bar u/t$&cal$_{\rm in}$&cal$^{\rm OOS}_{\rm vis}$&cal$^{\rm OOS}_{\rm sub}$\\\midrule
1 & 500 & 0.0000 & 0.0001 & 0.06 & 0.178 & 0.27 & 0.0017 & 0.0197 & 0.0220 \\
1 & 1000 & 0.0000 & 0.0000 & 0.01 & 0.057 & 0.05 & 0.0011 & 0.0160 & 0.0153 \\
1 & 2000 & 0.0000 & 0.0000 & 0.00 & 0.000 & -- & 0.0007 & 0.0127 & 0.0112 \\
2 & 500 & 0.0000 & 0.0003 & 0.15 & 0.652 & 0.34 & 0.0018 & 0.0194 & 0.0222 \\
2 & 1000 & 0.0000 & 0.0000 & 0.02 & 0.201 & 0.12 & 0.0011 & 0.0160 & 0.0148 \\
2 & 2000 & 0.0000 & 0.0000 & 0.00 & 0.000 & -- & 0.0007 & 0.0122 & 0.0114 \\
\bottomrule\end{tabular}\end{table}


\subsection{Onset-distribution robustness (non-proportional hazards / AFT) and pilot}
\begin{table}[t]\centering\small
\caption{Pilot results for the quantile slope $\beta_{\tau,1}$, \textbf{IPCW}
pipeline (IPCW--A-CV, IPCW--A-Rich, IPCW--B-OS-NCF). Bias is
Monte-Carlo over $R=200$ replicates; coverages are over a $150$-replicate
bootstrap subset (Monte-Carlo SE $\approx0.018$).
``boot'' = full-refit subject bootstrap; ``froz'' = frozen-$G_C$ bootstrap.
A separate paired CC-vs-IPCW study is in Table~\ref{tab:cc} (distinct
experiment; the IPCW values here are not reused there).}
\label{tab:pilot}
\begin{tabular}{rrr r rr r rrrr}
\toprule
Sc. & $n$ & $\tau$ & $\beta_{\tau,1}$ & bias$_{\text{ACV}}$ & bias$_{\text{BOS}}$ & ESD$_{\text{BOS}}$ & cov$_{\text{ACV}}^{\text{boot}}$ & cov$_{\text{ARich}}^{\text{boot}}$ & cov$_{\text{BOS}}^{\text{boot}}$ & cov$_{\text{BOS}}^{\text{froz}}$ \\
\midrule
1 & 500 & 0.10 & 0.500 & 0.017 & 0.018 & 0.145 & 0.907 & 0.927 & 0.900 & 0.900 \\
1 & 500 & 0.25 & 0.500 & 0.001 & 0.001 & 0.132 & 0.880 & 0.907 & 0.887 & 0.887 \\
1 & 500 & 0.50 & 0.500 & 0.008 & 0.007 & 0.106 & 0.953 & 0.967 & 0.960 & 0.960 \\
1 & 500 & 0.75 & 0.500 & -0.023 & -0.024 & 0.111 & 0.940 & 0.960 & 0.947 & 0.947 \\
1 & 500 & 0.90 & 0.500 & -0.005 & -0.004 & 0.139 & 0.933 & 0.940 & 0.953 & 0.953 \\
\addlinespace
1 & 1000 & 0.10 & 0.500 & 0.007 & 0.007 & 0.107 & 0.900 & 0.893 & 0.907 & 0.907 \\
1 & 1000 & 0.25 & 0.500 & 0.007 & 0.007 & 0.087 & 0.960 & 0.947 & 0.960 & 0.953 \\
1 & 1000 & 0.50 & 0.500 & -0.009 & -0.010 & 0.083 & 0.947 & 0.947 & 0.947 & 0.947 \\
1 & 1000 & 0.75 & 0.500 & 0.003 & 0.004 & 0.083 & 0.933 & 0.953 & 0.933 & 0.933 \\
1 & 1000 & 0.90 & 0.500 & -0.003 & -0.003 & 0.102 & 0.927 & 0.907 & 0.920 & 0.920 \\
\addlinespace
1 & 2000 & 0.10 & 0.500 & -0.002 & -0.002 & 0.071 & 0.940 & 0.960 & 0.947 & 0.947 \\
1 & 2000 & 0.25 & 0.500 & 0.005 & 0.005 & 0.059 & 0.960 & 0.947 & 0.960 & 0.960 \\
1 & 2000 & 0.50 & 0.500 & -0.000 & 0.000 & 0.057 & 0.927 & 0.933 & 0.933 & 0.933 \\
1 & 2000 & 0.75 & 0.500 & 0.007 & 0.007 & 0.053 & 0.953 & 0.987 & 0.967 & 0.967 \\
1 & 2000 & 0.90 & 0.500 & -0.006 & -0.005 & 0.071 & 0.953 & 0.933 & 0.953 & 0.953 \\
\addlinespace
2 & 500 & 0.10 & -0.013 & 0.027 & 0.028 & 0.176 & 0.913 & 0.900 & 0.913 & 0.913 \\
2 & 500 & 0.25 & 0.230 & -0.002 & -0.003 & 0.145 & 0.933 & 0.947 & 0.940 & 0.940 \\
2 & 500 & 0.50 & 0.500 & 0.005 & 0.006 & 0.144 & 0.953 & 0.953 & 0.960 & 0.960 \\
2 & 500 & 0.75 & 0.770 & -0.024 & -0.026 & 0.149 & 0.947 & 0.933 & 0.947 & 0.947 \\
2 & 500 & 0.90 & 1.013 & -0.017 & -0.017 & 0.182 & 0.933 & 0.933 & 0.933 & 0.927 \\
\addlinespace
2 & 1000 & 0.10 & -0.013 & -0.005 & -0.005 & 0.135 & 0.933 & 0.920 & 0.927 & 0.927 \\
2 & 1000 & 0.25 & 0.230 & 0.014 & 0.015 & 0.097 & 0.967 & 0.960 & 0.973 & 0.967 \\
2 & 1000 & 0.50 & 0.500 & -0.000 & 0.000 & 0.110 & 0.947 & 0.940 & 0.947 & 0.947 \\
2 & 1000 & 0.75 & 0.770 & -0.003 & -0.003 & 0.103 & 0.967 & 0.967 & 0.967 & 0.967 \\
2 & 1000 & 0.90 & 1.013 & -0.004 & -0.003 & 0.121 & 0.953 & 0.933 & 0.953 & 0.953 \\
\addlinespace
2 & 2000 & 0.10 & -0.013 & 0.000 & 0.000 & 0.084 & 0.960 & 0.960 & 0.967 & 0.967 \\
2 & 2000 & 0.25 & 0.230 & -0.006 & -0.006 & 0.073 & 0.953 & 0.960 & 0.953 & 0.953 \\
2 & 2000 & 0.50 & 0.500 & 0.005 & 0.005 & 0.076 & 0.927 & 0.920 & 0.927 & 0.927 \\
2 & 2000 & 0.75 & 0.770 & -0.007 & -0.007 & 0.075 & 0.940 & 0.940 & 0.933 & 0.933 \\
2 & 2000 & 0.90 & 1.013 & 0.009 & 0.009 & 0.091 & 0.960 & 0.947 & 0.947 & 0.947 \\
\bottomrule
\end{tabular}
\end{table}

\begin{table}[t]\centering\small
\caption{Non-proportional-hazards (log-normal AFT) onset, \textbf{IPCW--A-CV} pipeline (as in Table~\ref{tab:pilot}). Bias and bootstrap coverage of $\beta_{\tau,1}$ are broadly similar to the Cox-Weibull results of Table~\ref{tab:pilot}, supporting robustness to the onset mechanism (the method never models onset). $R=200$ for bias; coverage over a \emph{60}-replicate subset only (Monte-Carlo SE $\approx0.028$), so this is a \emph{pilot-grade} stress test --- the several $0.90$--$0.917$ cells are not distinguishable from nominal at this $R$. The $-0.050$ bias at (Sc.\ 2, $n{=}500$, $\tau{=}0.9$) has ESD $\approx0.18$ (bias MCSE $\approx0.023$), so it is $\approx2$ MCSE and not yet alarming; a CC--A-CV AFT run and a larger $R$ are deferred to the consolidated specification.}
\label{tab:nonph}\begin{tabular}{rrr r rr}\toprule
Sc.&$n$&$\tau$&$\beta_{\tau,1}$&bias$_{\text{ACV}}$&cov$_{\text{ACV}}^{\text{boot}}$\\\midrule
1 & 500 & 0.100 & 0.500 & 0.003 & 0.950
\\
1 & 500 & 0.500 & 0.500 & -0.003 & 0.933
\\
1 & 500 & 0.900 & 0.500 & -0.007 & 0.933
\\
\addlinespace
1 & 1000 & 0.100 & 0.500 & -0.005 & 0.917
\\
1 & 1000 & 0.500 & 0.500 & 0.008 & 0.917
\\
1 & 1000 & 0.900 & 0.500 & -0.009 & 0.983
\\
\addlinespace
1 & 2000 & 0.100 & 0.500 & 0.002 & 0.950
\\
1 & 2000 & 0.500 & 0.500 & -0.007 & 0.983
\\
1 & 2000 & 0.900 & 0.500 & -0.008 & 0.967
\\
\addlinespace
2 & 500 & 0.100 & -0.013 & 0.008 & 0.967
\\
2 & 500 & 0.500 & 0.500 & 0.008 & 0.933
\\
2 & 500 & 0.900 & 1.013 & -0.050 & 0.900
\\
\addlinespace
2 & 1000 & 0.100 & -0.013 & 0.005 & 0.967
\\
2 & 1000 & 0.500 & 0.500 & -0.016 & 0.983
\\
2 & 1000 & 0.900 & 1.013 & -0.018 & 0.950
\\
\addlinespace
2 & 2000 & 0.100 & -0.013 & 0.011 & 0.900
\\
2 & 2000 & 0.500 & 0.500 & 0.005 & 0.950
\\
2 & 2000 & 0.900 & 1.013 & -0.005 & 0.900
\\
\bottomrule\end{tabular}\end{table}


\subsection{Computational reproducibility audit}
\begin{table}[t]\centering\small
\caption{Computational reproducibility audit, a \emph{dedicated reduced-$B$ run} ($R=300$ per cell, inner $B{=}60$; \emph{not} a summary of the headline $B{=}500$ log), $\tau\in\{.25,.5,.75\}$. Failure/fallback rates that would bias the headline numbers if silently dropped, for \emph{both} the CC and IPCW pipelines: the QR (QR-solver) and CF (B-OS-CF) failure columns are shown \emph{separately for both pipelines}; the remaining three columns---$\widehat f$-cap, $\kappa(\widehat S)$, and the inner subject-bootstrap replicate-failure rate (``boot'')---are \emph{CC-pipeline} diagnostics (the recommended estimator). \emph{Failed replicates are counted, not discarded}. We also report the bread condition-number distribution (median, $95$th pct) rather than only a permissive (and loose) $\kappa>10^8$ binary flag. \emph{Denominators}: QR and CF are over the $R{=}300$ datasets; $\widehat f$-cap over all fold--visit density estimates; boot over the $R{\times}B$ inner resamples. \textbf{The QR and B-OS-CF structural-fit failure rates are $0$ for both the CC and IPCW pipelines}, and the (CC) bootstrap-replicate failure rate is $0$ (we tabulate the CC bootstrap-failure column only, so we do \emph{not} claim an IPCW bootstrap-failure audit here); a zero-event rate at $R{=}300$ has a $95\%$ binomial upper bound of $\approx1\%$, not an exact $0$. $\kappa(\widehat S)$ stays small (median $\approx4$, $95$th pct $\le9.1$), so ill-conditioning is not a concern; the only non-zero mode is the $\widehat f$-cap ($\le1.0\%$). \emph{Scope}: base CC/IPCW pipelines at the central quartile quantiles only; the degenerate small-$n$ growing-sieve bootstrap (Table~\ref{tab:sieve}) is \emph{not} covered here and needs a separate distributional audit.}
\label{tab:audit}\begin{tabular}{rrr rr rr rr rr}\toprule
&&&\multicolumn{2}{c}{CC}&\multicolumn{2}{c}{IPCW}&\multicolumn{2}{c}{$\kappa(\widehat S)$}&&\\
\cmidrule(lr){4-5}\cmidrule(lr){6-7}\cmidrule(lr){8-9}
Sc.&$n$&$\tau$&QR&CF&QR&CF&med&$q_{95}$&$\widehat f$-cap&boot\\\midrule
1 & 500 & 0.25 & 0.000 & 0.000 & 0.000 & 0.000 & 4.0 & 6.5 & 0.009 & 0.000 \\
1 & 500 & 0.50 & 0.000 & 0.000 & 0.000 & 0.000 & 4.1 & 6.1 & 0.009 & 0.000 \\
1 & 500 & 0.75 & 0.000 & 0.000 & 0.000 & 0.000 & 4.1 & 6.5 & 0.008 & 0.000 \\
\addlinespace
1 & 1000 & 0.25 & 0.000 & 0.000 & 0.000 & 0.000 & 4.0 & 5.2 & 0.003 & 0.000 \\
1 & 1000 & 0.50 & 0.000 & 0.000 & 0.000 & 0.000 & 3.9 & 5.0 & 0.003 & 0.000 \\
1 & 1000 & 0.75 & 0.000 & 0.000 & 0.000 & 0.000 & 3.9 & 5.1 & 0.003 & 0.000 \\
\addlinespace
1 & 2000 & 0.25 & 0.000 & 0.000 & 0.000 & 0.000 & 3.9 & 4.6 & 0.001 & 0.000 \\
1 & 2000 & 0.50 & 0.000 & 0.000 & 0.000 & 0.000 & 3.8 & 4.4 & 0.001 & 0.000 \\
1 & 2000 & 0.75 & 0.000 & 0.000 & 0.000 & 0.000 & 3.8 & 4.5 & 0.001 & 0.000 \\
\addlinespace
2 & 500 & 0.25 & 0.000 & 0.000 & 0.000 & 0.000 & 4.9 & 9.0 & 0.010 & 0.000 \\
2 & 500 & 0.50 & 0.000 & 0.000 & 0.000 & 0.000 & 4.9 & 8.6 & 0.009 & 0.000 \\
2 & 500 & 0.75 & 0.000 & 0.000 & 0.000 & 0.000 & 4.9 & 9.1 & 0.010 & 0.000 \\
\addlinespace
2 & 1000 & 0.25 & 0.000 & 0.000 & 0.000 & 0.000 & 4.4 & 6.5 & 0.003 & 0.000 \\
2 & 1000 & 0.50 & 0.000 & 0.000 & 0.000 & 0.000 & 4.3 & 6.2 & 0.003 & 0.000 \\
2 & 1000 & 0.75 & 0.000 & 0.000 & 0.000 & 0.000 & 4.4 & 6.9 & 0.004 & 0.000 \\
\addlinespace
2 & 2000 & 0.25 & 0.000 & 0.000 & 0.000 & 0.000 & 4.2 & 5.1 & 0.001 & 0.000 \\
2 & 2000 & 0.50 & 0.000 & 0.000 & 0.000 & 0.000 & 4.2 & 5.2 & 0.001 & 0.000 \\
2 & 2000 & 0.75 & 0.000 & 0.000 & 0.000 & 0.000 & 4.2 & 4.9 & 0.001 & 0.000 \\
\bottomrule\end{tabular}\end{table}


\subsection{Comparison with the conditional-mean estimator of Sun et al.}
\label{supp:sun}
Full numbers for the five-regime comparison of Section~\ref{sec:suncmp}.
Table~\ref{tab:cmp-beta} reports distributional-effect recovery and scalar-slope
variability (with the oracle-$\alpha$ decomposition discussed in the main text);
Table~\ref{tab:cmp-contrast} the mean-contrast IMSE for Sun-basic, the stratified
Sun-S extension, and the two quantile reconstructions Q-LN and Q-INT, with paired
Monte-Carlo standard errors; Table~\ref{tab:cmp-surface} the conditional-mean
surface and Sun bandwidth sensitivity; and Table~\ref{tab:cmp-heavy} the
heavy-tail (DGP5) robustness diagnostics.
\begin{table}[t]\centering\small
\caption{\textbf{Distributional effect recovery and scalar-slope variability (Sun's real estimator vs.\ ours).} Same Monte-Carlo datasets under Sun et al.'s survival/truncation/fixed-visit design ($R=500$; common interior domain). Sun's profile partial-likelihood slope $\widehat\beta_1$ (their compiled \texttt{marker3.cpp}) and its Monte-Carlo ESD vs.\ the ESD of our median slope $\widehat\beta_{.5,1}$, plus our $\beta_{\tau,1}$ bias at $\tau=0.1,0.5,0.9$ and the recovered tail spread $\widehat\beta_{.9}-\widehat\beta_{.1}$. \emph{Both estimators are unbiased when their target is well-defined.} In DGP1 the direct median slope has smaller Monte-Carlo variability than the feasible Sun slope for the shared scalar target ($\sim$$4\times$ smaller ESD); this is a point-estimation comparison, not an inferential-efficiency claim. An oracle-$\alpha$ decomposition shows Sun's feasible $\widehat\beta=\widehat\theta-\widehat\alpha$ (Thm.~3.1) is \emph{not} less stable because it estimates $\alpha$: the oracle-$\alpha$ version $\widehat\theta-\alpha_0$ is in fact \emph{more} variable, so the dominant variability appears to come from the profile-kernel marker score $\widehat\theta$, with $\widehat\theta-\widehat\alpha$ partly benefiting from covariance cancellation. \emph{The Sun-ESD vs.\ Q-ESD comparison is like-for-like only in DGP1, where the mean and median slopes coincide ($\beta_{\rm mean}=\beta_{.5}$); in DGP2--DGP4 the scalar targets differ, and in DGP5 the conditional mean is undefined, so Sun's slope variability is reported only as a fragility diagnostic, not an efficiency comparison.} DGP2 only the quantile slope sees the covariate-driven dispersion; DGP3 no proportional slope exists; DGP4 the effect is confined to the upper tail ($\beta_{.9}\!\gg\!\beta_{.5}$) so Sun's scalar hides it; in DGP5 (heavy $t_5$ log-tail, mean ill-posed) Sun's raw ESD is unstable and non-monotone, driven by rare explosive observations, while robust spread summaries shrink with $n$ (Table~\ref{tab:cmp-heavy}) and our median slope is robust.}
\label{tab:cmp-beta}\begin{tabular}{r rr rrr r}\toprule
&&&\multicolumn{3}{c}{Q $\beta_{\tau,1}$ bias}&spread $\beta_{.9}{-}\beta_{.1}$\\
\cmidrule(lr){4-6}
$n$&Sun $\widehat\beta_1$ (ESD)&Q ESD$_{.5}$&$.10$&$.50$&$.90$&Q / true\\\midrule
\multicolumn{7}{l}{\textbf{DGP1} proportional, homoscedastic mean}\\
500 & 0.51\,(0.10) & 0.024 & -0.002 & -0.001 & 0.002 & 0.00/0.00\\
1000 & 0.50\,(0.07) & 0.017 & 0.001 & 0.000 & 0.001 & 0.00/0.00\\
2000 & 0.50\,(0.05) & 0.012 & 0.001 & 0.001 & 0.001 & 0.00/0.00\\
\addlinespace
\multicolumn{7}{l}{\textbf{DGP2} proportional mean, dispersion varies}\\
500 & 0.50\,(0.15) & 0.088 & 0.008 & -0.003 & -0.019 & 1.13/1.15\\
1000 & 0.50\,(0.11) & 0.063 & 0.010 & 0.001 & -0.001 & 1.14/1.15\\
2000 & 0.50\,(0.08) & 0.047 & 0.000 & -0.006 & -0.006 & 1.15/1.15\\
\addlinespace
\multicolumn{7}{l}{\textbf{DGP3} non-proportional mean (headline)}\\
500 & 0.77\,(0.14) & 0.063 & 0.013 & -0.001 & -0.014 & 1.13/1.15\\
1000 & 0.76\,(0.10) & 0.046 & 0.008 & 0.001 & -0.006 & 1.14/1.15\\
2000 & 0.77\,(0.07) & 0.032 & 0.007 & 0.001 & -0.002 & 1.14/1.15\\
\addlinespace
\multicolumn{7}{l}{\textbf{DGP4} tail-specific effect}\\
500 & 0.67\,(0.20) & 0.064 & 0.004 & 0.005 & -0.111 & 0.98/1.10\\
1000 & 0.67\,(0.14) & 0.047 & -0.002 & 0.003 & -0.051 & 1.05/1.10\\
2000 & 0.67\,(0.09) & 0.033 & 0.000 & -0.001 & -0.031 & 1.07/1.10\\
\addlinespace
\multicolumn{7}{l}{\textbf{DGP5} heavy-tailed error (mean undefined)}\\
500 & 0.53\,(0.49) & 0.054 & 0.000 & 0.001 & 0.002 & 0.00/0.00\\
1000 & 0.50\,(0.09) & 0.037 & 0.002 & 0.002 & 0.002 & 0.00/0.00\\
2000 & 0.52\,(0.37) & 0.028 & -0.002 & -0.001 & 0.000 & 0.00/0.00\\
\bottomrule\end{tabular}\end{table}

\begin{table}[t]\centering\small
\caption{\textbf{Conditional-mean contrast --- the headline} ($\mathrm{IMSE}\times10^{3}$). Both methods target the conditional mean. Sun: the \emph{basic} proportional model ($E[Y|u,t,x]=\mu_0(u,t)e^{\beta'x}$, log-contrast $=\widehat\beta_1$, constant) and his Sec.~3.5 \emph{stratified-baseline} extension (Sun-S: his genuine estimator fit separately per $X_1$ group, giving a nonparametric $(u,t)$-varying contrast). Ours: the model-free quantile integral Q-INT ($\int_0^1 e^{Q_\tau}d\tau$) and the tail-stable retransformation Q-LN ($\widehat Q_{.5}+\tfrac12\widehat\sigma^2$). IMSE of the mean log-contrast $\Delta(u,t)$. Under non-proportionality (DGP3) $\Delta$ truly varies: Sun-basic collapses it to a pseudo-average, and the fully stratified Sun-S recovers it only at a large \emph{variance} cost (two nonparametric surfaces), so its IMSE exceeds Sun-basic and far exceeds Q (both Q reconstructions several-fold lower; the paired $\text{Q-LN}-\text{Sun-basic}$ difference is significantly negative, MCSE in parentheses). Sun-S is a fully stratified implementation and is \emph{not} intended to represent all prespecified low-dimensional interaction bases. In DGP1 the lower Q IMSE reflects the smaller Monte-Carlo variability of the direct shared scalar slope (Table~\ref{tab:cmp-beta}); in DGP2 Q-LN benefits from the log-location-scale structure, whereas the model-free Q-INT is more tail-sensitive and is \emph{not} uniformly better than Sun ($\text{Q-INT}-\text{Sun}=+2.5\times10^{-3}$ at $n{=}2000$). In \textbf{DGP4 Sun-basic is \emph{best}} (well matched to the proportional mean target). DGP5's mean is undefined ($-$).}
\label{tab:cmp-contrast}\begin{tabular}{r rrrr rr}\toprule
&\multicolumn{4}{c}{contrast IMSE $\times10^{3}$}&\multicolumn{2}{c}{paired $-$Sun $\times10^{3}$ (MCSE)}\\
\cmidrule(lr){2-5}\cmidrule(lr){6-7}
$n$&Sun&Sun-S&Q-LN&Q-INT&Q-LN&Q-INT\\\midrule
\multicolumn{7}{l}{\textbf{DGP1} proportional, homoscedastic mean}\\
500 & 9.445 & 93.959 & 0.360 & 0.353 & -9.084\,(0.624) & -9.092\,(0.624)\\
1000 & 5.084 & 42.660 & 0.183 & 0.184 & -4.901\,(0.312) & -4.900\,(0.312)\\
2000 & 2.329 & 20.742 & 0.091 & 0.091 & -2.238\,(0.151) & -2.239\,(0.151)\\
\addlinespace
\multicolumn{7}{l}{\textbf{DGP2} proportional mean, dispersion varies}\\
500 & 22.763 & 188.000 & 8.334 & 14.944 & -14.429\,(1.434) & -7.820\,(1.558)\\
1000 & 12.612 & 102.526 & 4.546 & 9.597 & -8.066\,(0.710) & -3.015\,(0.797)\\
2000 & 5.853 & 59.003 & 2.306 & 8.374 & -3.547\,(0.350) & 2.522\,(0.452)\\
\addlinespace
\multicolumn{7}{l}{\textbf{DGP3} non-proportional mean (headline)}\\
500 & 23.703 & 158.101 & 4.398 & 7.782 & -19.305\,(1.275) & -15.921\,(1.338)\\
1000 & 14.290 & 78.095 & 2.169 & 5.104 & -12.121\,(0.580) & -9.187\,(0.634)\\
2000 & 9.539 & 41.998 & 1.043 & 3.876 & -8.496\,(0.320) & -5.663\,(0.373)\\
\addlinespace
\multicolumn{7}{l}{\textbf{DGP4} tail-specific effect}\\
500 & 40.235 & 239.699 & 53.970 & 27.912 & 13.734\,(3.103) & -12.323\,(2.463)\\
1000 & 18.489 & 139.674 & 40.737 & 14.536 & 22.247\,(1.794) & -3.953\,(1.210)\\
2000 & 8.761 & 72.335 & 36.858 & 9.927 & 28.097\,(1.239) & 1.166\,(0.695)\\
\addlinespace
\multicolumn{7}{l}{\textbf{DGP5} heavy-tailed error (mean undefined)}\\
500 & -- & -- & -- & -- & --\,(--) & --\,(--)\\
1000 & -- & -- & -- & -- & --\,(--) & --\,(--)\\
2000 & -- & -- & -- & -- & --\,(--) & --\,(--)\\
\bottomrule\end{tabular}\end{table}

\begin{table}[t]\centering\small
\caption{\textbf{Conditional-mean surface and Sun bandwidth sensitivity} (log-IMSE $\times10^{3}$, $x_1{=}0$). Sun's kernel baseline $\widehat\mu_0$ at his density-reference bandwidth $h$ ($2.34\,\mathrm{sd}(A)N^{-1/6}$) and at the regression bandwidth $0.5h$, vs.\ our Q-LN and Q-INT reconstructions. Sun's surface is bandwidth-sensitive---the density rule over-smooths the steep baseline ($h$ column), and $0.5h$ roughly halves the IMSE---so the Sun-vs-Q surface gap is partly a smoothing-bandwidth effect, not a method deficiency; the bandwidth-free slope and contrast (Tables~\ref{tab:cmp-beta}--\ref{tab:cmp-contrast}) are the substantive comparisons. DGP5's mean is undefined ($-$).}
\label{tab:cmp-surface}\begin{tabular}{rr rr rr}\toprule
&&\multicolumn{2}{c}{Sun $\widehat\mu_0$}&\multicolumn{2}{c}{ours}\\
\cmidrule(lr){3-4}\cmidrule(lr){5-6}
$n$&$R$&$h$&$0.5h$&Q-LN&Q-INT\\\midrule
\multicolumn{6}{l}{\textbf{DGP1} proportional, homoscedastic mean}\\
500 & 500 & 31.021 & 15.380 & 1.081 & 0.967\\
1000 & 500 & 18.505 & 8.316 & 0.695 & 0.607\\
2000 & 500 & 10.757 & 4.509 & 0.519 & 0.442\\
\addlinespace
\multicolumn{6}{l}{\textbf{DGP2} proportional mean, dispersion varies}\\
500 & 500 & 36.480 & 29.638 & 5.053 & 5.314\\
1000 & 500 & 22.935 & 18.114 & 2.625 & 2.948\\
2000 & 500 & 12.885 & 10.055 & 1.428 & 1.826\\
\addlinespace
\multicolumn{6}{l}{\textbf{DGP3} non-proportional mean (headline)}\\
500 & 500 & 40.442 & 30.410 & 2.400 & 2.411\\
1000 & 500 & 24.352 & 17.015 & 1.217 & 1.293\\
2000 & 500 & 15.570 & 10.565 & 0.754 & 0.848\\
\addlinespace
\multicolumn{6}{l}{\textbf{DGP4} tail-specific effect}\\
500 & 500 & 43.106 & 39.571 & 4.737 & 4.582\\
1000 & 500 & 26.898 & 24.246 & 2.388 & 2.566\\
2000 & 500 & 15.794 & 13.366 & 1.305 & 1.482\\
\addlinespace
\multicolumn{6}{l}{\textbf{DGP5} heavy-tailed error (mean undefined)}\\
500 & 500 & -- & -- & -- & --\\
1000 & 500 & -- & -- & -- & --\\
2000 & 500 & -- & -- & -- & --\\
\bottomrule\end{tabular}\end{table}

\begin{table}[t]\centering\small
\caption{\textbf{DGP5 heavy-tail robustness diagnostics} (Student-$t_5$ log-error; true location shift $\beta_1=0.5$). Sun's mean-slope sampling distribution over $R=500$ replicates vs.\ our median slope. The ordinary ESD is large and non-monotone in $n$, but the median, IQR, $95$th percentile, and $3$-IQR-trimmed ESD (tESD) are small and \emph{shrink} with $n$: the inflated ESD is driven by a few replicates with an extreme maximum marker (the $t_5$ log-tail produces markers exceeding $10^4$ in rare replicates), not a coding artifact. Our median slope is unbiased with a small, shrinking ESD throughout. This is an estimand-validity illustration: the conditional mean (Sun's target) is undefined under a $t_5$ log-tail, so it is not a fair mean-estimation contest.}
\label{tab:cmp-heavy}\begin{tabular}{r rr rrrr r}\toprule
&\multicolumn{2}{c}{Sun $\widehat\beta_1$}&\multicolumn{4}{c}{Sun $\widehat\beta_1$ robust spread}&\\
\cmidrule(lr){2-3}\cmidrule(lr){4-7}
$n$&mean&ESD&med&IQR&$q_{.95}$&tESD&Q ESD$_{.5}$\\\midrule
500 & 0.53 & 0.49 & 0.51 & 0.15 & 0.69 & 0.12 & 0.054\\
1000 & 0.50 & 0.09 & 0.50 & 0.11 & 0.63 & 0.08 & 0.037\\
2000 & 0.52 & 0.37 & 0.50 & 0.08 & 0.60 & 0.06 & 0.028\\
\bottomrule\end{tabular}\end{table}


% ============================================================================
\section{Additional application details}\label{supp:app}

\paragraph{Synthetic cohort.} The illustrative analysis of Section~\ref{sec:app}
uses a synthetic cohort generated to match the ADNI sampling design: covariates are
a binary \emph{APOE4} carrier indicator and sex; onset (amyloid positivity) is
Cox--Weibull; competing death is Weibull; entry is staggered with left truncation;
loss to follow-up is independent given $(A,\X)$; visits are noninformative. The
log-marker is generated with a covariate effect that varies across quantile levels
so that the carrier effect $\beta_{\tau,1}$ is non-flat in $\tau$, illustrating a
distributional (not pure-location) effect. \emph{These are not real ADNI data and
the figures are for illustration of the deliverables only; the substantive ADNI
analysis is in preparation.}

\paragraph{Estimable region and masks.} Quantile surfaces are reported on the
interior region where the visit occupation density is bounded away from zero; grid
cells with fewer than a minimum number of nearby complete-event visits are masked.

\paragraph{Sensitivity.} For the real-data analysis we will report sensitivity to
(i) the weighting (CC default vs.\ IPCW design-standardization), (ii) the
inferential route (A-CV with subject bootstrap vs.\ B-OS-CF), and (iii) crossing
(rearrangement of predicted curves), following the same protocol validated in
Section~\ref{supp:sims}.

\end{document}
